%% Version 1.2.17
\documentclass[12pt,letterpaper]{book}
\usepackage{amsmath,amssymb,epic,eepic,pstricks,xcolor,comment}
\usepackage[body={6in,8.5in}]{geometry}
%\usepackage{makeidx}
\usepackage{index}
%\newindex{xfcns}{fdx}{fnd}{Index of Functions}
\makeindex

\usepackage{ifpdf}

\usepackage[
  \ifpdf pdftex, \else dvips, \fi
  hyperfootnotes=false,
  pdftitle={ePiX Tutorial and Reference Manual},
  pdfauthor={Andrew D. Hwang},
  pdfkeywords={LaTeX, Mathematical Figures, Plotting},
  pdfstartview=FitH,
  bookmarks=true,
  colorlinks=true,
  linkcolor=black]{hyperref}

\setlength{\emergencystretch}{1.25em}
\renewcommand{\thefigure}{\arabic{chapter}.\arabic{figure}}

\newcommand{\code}[1]{\texttt{#1}}
\newcommand{\filename}[1]{\texttt{#1}}
\newcommand{\Key}[1]{\texttt{#1}}
\newcommand{\URL}[1]{\texttt{#1}}

%% \setlength{\topmargin}{0.5in}
\newcommand{\ePiX}{\code{ePiX}}
\newcommand{\ext}{\code{ePiX\_ext}}
\newcommand{\epix}{\code{epix}}
\newcommand{\eepic}{\code{eepic}}
\newcommand{\C}{\code{C}}
\newcommand{\CXX}{\code{C++}}
\newcommand{\pyepix}{\code{Pyepix}}

\DeclareMathOperator{\clip}{clip}

\newcommand{\Header}[1]{\smallskip\noindent\texttt{#1}\quad}%

%% Substitution commands easily removable by sed for texinfo:
%% Remove recalcitrant chunks,
\newcommand{\texicomment}{}
\newcommand{\texicommentend}{}

%% mark simple math,
\newcommand{\teximath}[1]{$#1$}
%% make function index entries,
%\newcommand{\findex}[1]{\index[xfcns]{#1}}
\newcommand{\findex}[1]{\ignorespaces}

%% top-level division (first argument for texinfo node name)
\newcommand{\texinfochapter}[2]{\chapter{#2}}
\newcommand{\texinfoappendix}[2]{\chapter{#2}}
%%\renewcommand{\subsection*}[1]{\subsection{#1}}

\title{ \ePiX\ Tutorial and  \\ Reference Manual}
\author{Andrew D. Hwang \\ 
  Department\ of Math and CS \\ 
  College of the Holy Cross}
\date{Version 1.2.17, July, 2017}

\begin{document}

\maketitle

\tableofcontents

\texinfochapter{1}{Introduction}

\ePiX, a collection of batch utilities, creates mathematically
accurate figures, plots, and animations containing \LaTeX\ typography.
The input syntax is easy to learn, and the user interface resembles
that of \LaTeX\ itself: You prepare a scene description in a text
editor, then ``compile'' the input file into a picture. \LaTeX- and
web-compatible output types include a \LaTeX\ picture-like environment
written with \code{PSTricks}, \code{tikz}, or \code{eepic} macros;
vector images (\code{eps}, \code{ps}, and~\code{pdf}); and bitmapped
images and movies (\code{png}, \code{mng}, and~\code{gif}).

\ePiX's strengths include:
 
\begin{itemize}
\item Quality of output: \ePiX\ creates accurate, publication-quality
  figures whose appearance matches that of \LaTeX.  Typography may be
  put in a figure as easily as in an ordinary \LaTeX\ document.
  
\item Ease of use: Figure objects and their attributes are specified
  by simple, descriptive commands.
  
\item Flexibility: Objects are described by attributes and Cartesian
  location; as in \LaTeX, printed appearance is determined when the
  figure is compiled. A well-designed figure can be altered
  dramatically, yet precisely, with command-line switches or minor
  changes to the input file.
  
\item Power and extendibility: \ePiX\ inherits the power of \CXX\ as a
  programming language; variables, data structures, loops, and
  recursion can be used to draw complicated plots and figures with
  just a few lines of input. External code can be incorporated in a
  figure with a command line option or by using a Makefile.
  
\item Economy of storage and transmission: For a document containing
  many figures, a compressed tar file of the \LaTeX\ sources and
  \ePiX\ files is typically a few percent the size of the compressed
  PostScript file.
  
\item License: \ePiX\ is \emph{free software}.
  \index{Free software}
  You are granted the right to use the program for whatever purpose,
  and to inspect, modify, and re-distribute the source code, so long
  as you do not restrict the rights of others to do the same. In
  short, the license is similar to the terms under which theorems are
  published.

\end{itemize}


\ePiX\ facilitates logical structuring of mathematical figures, as
opposed to visual structuring (or WYSIWYG), analogous to the
relationship between \LaTeX\ and a word processor. Stylistic
defaults streamline the creation of simple figures, but there are few
internal restrictions on the contents or appearance of a figure;
aesthetic and practical decisions are left to you.

If you are a:

\begin{itemize}

\item Potential user, you may wish to skip immediately to ``Software
  Dependencies'' before investing additional time.

\item New user, proceed from here until you have enough understanding
  to run the software, then experiment with the samples files while
  reading Chapter~\ref{chapter-started}, or return to the manual as
  needed.

\item More advanced user, browse at will, probably starting with
  Chapter~\ref{chapter-ref-man}.

\end{itemize}

This manual is relatively conversational, and occasionally redundant,
especially between portions meant for readers at different levels of
familiarity. Throughout, you are assumed to be familiar with \LaTeX\
and basic linear algebra: the description of points, vectors, lines,
and planes in three-dimensional space. Other material, such as \CXX\
syntax, is introduced as needed.


\section{Software Dependencies}
\label{section-software_deps}

\index{Installation|(}
\index{Mac OS X}

If you run GNU/Linux, a BSD, or Solaris, you probably have (and can
surely install) all the external software needed to use \ePiX.  On Mac
OS~X, you will need the Apple developer tools and an X~server (such as
XQuartz), and the free~\code{fink} package manager to build a GNU
environment.  For Windows, you'll need to install Cygwin and several
packages. Detailed instructions are given below.

``Under the hood'', an input file is successively converted to a
\LaTeX\ picture; \code{dvi}; PostScript, \code{pdf}~or \code{eps}; and
if desired, to a bitmapped image or movie.  Four shell
scripts---\epix, \code{laps}, \code{elaps}, and~\code{flix}---automate
the various file format conversions.

\ePiX\ consists of a \CXX\ library, header, and shell scripts, and
requires GNU \code{bash} and a compiler \emph{for normal use}.  For
complete functionality, you need \code{g++} (Version~3.2 or later),
\code{bash}, a text editor (\ePiX\ works particularly well with
\code{emacs}), a \LaTeX\ distribution, \code{Ghostscript}, \code{gv}
(or your favorite PS/PDF previewer), and \code{GraphicsMagick}. GNU
\code{grep} and \code{sed} are good to have.  You may need additional
``developer packages'' (\code{binutils}, \code{make}) in order to
build \ePiX.  The more up to date your software is, the better your
experience is likely to be, but bleeding edge versions are not
necessary, or even always desirable.

Aside from their reliance on specific programs, \ePiX's shell scripts
are written using Unix-style pathnames. Thus, the most straightforward
way to use \ePiX\ is to install a GNU environment.

Jay Belanger's \code{emacs} mode allows you to write, compile, and
view \ePiX\ figures without leaving \code{emacs}. If you use another
editor, you'll want to create template source files so you don't have
to type boilerplate code each time you write a new figure.
\index{Emacs@\code{emacs}}


\subsection{Setting up an Environment Under  Windows}

\index{Windows operating system|(}
\begin{comment}
Version~1.0.4 of \ePiX\ has been implemented in Python~2.2 by Andrew
Sterian, making \ePiX\ available on any platform that supports Python,
and without requiring a \CXX~compiler or~\code{bash}.  Python is a
GPL-ed scripting language, and is available with a Windows installer
and detailed instructions.  The easiest alternative for Windows users
is probably to install Python~2.2 or later (if necessary) and
\code{Pyepix}. The \code{pyepix} project home page is:
\URL{claymore.engineer.gvsu.edu/~steriana/Python/index.html}
\end{comment}

Cygwin can be used to run \ePiX\ under Windows. Download
\filename{setup.exe} from \URL{www.cygwin.com}, then install the
packages you need. The following are recommended, and sufficient for
the actions described in this manual.
\begin{verbatim}
(Archive)    bzip2, tar
(Devel)      binutils, coreutils, gcc, gcc-g++, make, sed
(Editors)    emacs, emacs-X11, vim
(Graphics)   GraphicsMagick, ghostscript-base, ghostscript-x11, gv
(Publishing) tetex (all)
(Shells)     bash, bash-completion
(X11)        X-start-menu-icons, X-startup-scripts,
             XFree86-lib-compat, xorg-x11-fscl, xorg-x11-fsrv
\end{verbatim}
\index{Windows operating system|)}


\section{Installation}

\ePiX\ is distributed over the World-Wide Web as source code. Packages
may be found at \URL{mathcs.holycross.edu/\~{}ahwang/epix/ePiX.html}
The latest stable release is also on the CTAN mirrors, in the
\code{graphics} directory.  (Some users of Red Hat have reported file
permission problems when unpacking the CTAN tarballs. If you encounter
this difficulty, please try downloading the sources from the project
main page.)  Unpack the compressed tar file with the appropriate
command:
\begin{verbatim}
  tar -zxvf epix-x.y.z.tar.gz
  tar -jxvf epix-x.y.z.tar.bz2
\end{verbatim}
(\code{x.y.z}~is the version number) or, if your \code{tar}
doesn't do decompression,
\begin{verbatim}
  gunzip -c epix-x.y.z.tar.gz | tar -xvf -
  bzcat epix-x.y.z.tar.bz2 | tar -xvf -
\end{verbatim}
\code{cd} to the source directory, \filename{epix-x.y.z}.  The
\filename{INSTALL} file contains detailed installation instructions.
If you're impatient, the short of it is \code{./configure [--options];
  make; make install}.  Run \code{./configure --help} for a list of
options.

By default, \ePiX\ installs in subdirectories of
\filename{/usr/local}; if you want to install elsewhere, supply
\code{./configure} with the appropriate \code{--prefix}. You may also
want to consult \filename{POST-INSTALL} for information on setting
your \code{PATH} variable so your shell can find \ePiX.  The manual
and sample files are in \filename{/usr/local/share/doc/epix}.
\index{Installation|)}


\subsection{Development}

There are two mailing lists, one for user questions, one for
development discussion. Please visit
\URL{savannah.nongnu.org/mail/?group=epix} to subscribe.


\texinfochapter{2}{Getting Started}
\label{chapter-started}

This chapter describes the basics of creating figures in \ePiX\ for
readers familiar with \LaTeX\ but completely new to~\CXX. No detailed
knowledge of~\CXX\ is needed to use \ePiX, only a bit of grammar that
is easily absorbed by example.

Section~\ref{section-running} describes the commands (shell scripts)
comprising \ePiX, and explains how to set up a graphical environment
using standard *nix programs. Section~\ref{section-overview} briefly
describes figure creation. Section~\ref{section-tutorial} presents a
few files side-by-side with their output, and should be read at a
computer so you can run the exercises.


\section{Running \ePiX}
\label{section-running}

An ``input file'' is a human-written figure specification containing
\ePiX\ commands; an ``output file'' is either a \LaTeX\ picture-like
environment or a stand-alone vector graphic.  Conversion
(``compiling'' a figure) is accomplished with four shell scripts,
\code{laps}, \code{epix}, \code{elaps}, and \code{flix}.  Each script
has a preferred extension for its input files, and is invoked with a
command of the form
\begin{verbatim}
  <script> [options] <input file(s)>
\end{verbatim}
Often, no options are necessary.  \code{<script> --help}
describes \code{<script>}'s options.
\texicomment
Figure~\ref{fig-programs} (page~\pageref{fig-programs}) diagrams the
shell scripts and the file types they process.
\texicommentend

By default, output file names are constructed by replacing the input
extension with the (preferred) output extension. For brevity,
extensions may be omitted. If the script has doubts about your intent,
it proceeds with default behavior and prints a warning message.

The author is a great fan of \Key{TAB} completion, under which a
shell, based on what has been typed so far, automatically fills in a
command when the \Key{TAB} key is pressed.  \ePiX\ comes with code
snippets that complete intelligently when the first part of a command
is one of the shell scripts. For example, if command completion is
active, typing \code{epix} \Key{TAB} prints only names of \code{epix}
input files. To use this feature, you must install Ian MacDonald's
\code{bash completion} package. The \filename{INSTALL} and
\filename{POST\_INSTALL} files contain details.


\subsection*{Shell Scripts}

\subsubsection*{\code{laps}}

\code{laps} performs \LaTeX\ to PostScript/PDF conversion, and is
independent of the rest of \ePiX. By default, \code{laps} invokes
\LaTeX~and \code{dvips}. The option \code{--pdf} creates a PDF~file by
post-processing the PostScript with~\code{ps2pdf}.  Other \TeX-family
processors (\code{pslatex}, \code{pdftex}, etc.) may be used instead
of~\LaTeX\ by invoking \code{laps} with an appropriate option.


\subsubsection*{\code{epix}}

\code{epix} compiles an input file into a \LaTeX\ picture.  The input
file should have extension \filename{xp} (for eXtended Picture).  Jay
Belanger's \code{emacs} mode is Highly Recommended; the installation
procedure can be found in the \filename{POST\_INSTALL} file. When this
mode is active, \code{emacs} automatically inserts a preamble template
when a new \filename{xp} file is created. In addition, you can peruse
\ePiX's info file, or format, compile, and preview files, all from
within \code{emacs}.
\index{Emacs@\code{emacs}}

Output files from the script \code{epix} have extension
\filename{eepic}, after Conrad Kwok's enhancements to the \LaTeX\
picture environment. In Version~1.2, however, an output file may
contain \code{PSTricks} or \code{tikz} macros instead of \code{eepic}
macros.  A command such as \verb+\input{myfile.eepic}+ incorporates an
output file into a \LaTeX\ document. The preamble must contain
\code{usepackage} lines appropriate to the output format.
\texicomment
\begin{center}
  \begin{tabular}{cc}
    Format & Required Package(s) \\
    \hline
    \code{eepic} & \code{epic,eepic,xcolor,rotating} \\
    \code{pst}   & \code{pstricks,rotating} \\
    \code{tikz}  & \code{tikz,rotating}
  \end{tabular}
\end{center}
\texicommentend


\subsubsection*{\code{elaps}}

\code{elaps} creates stand-alone vector images (\filename{eps}~or
\filename{pdf}) from \code{epix} input files or \filename{eepic}
figures, even those not produced with \code{epix}.  \code{elaps}
automatically loads the \LaTeX\ packages needed for features described
in this manual. Additional \LaTeX\ packages and \code{dvips} options
may be specified on the command line.  An \code{elaps} output file is
placed in a \LaTeX\ document with an \verb+\includegraphics+ command;
the preamble must contain \verb+\usepackage{graphicx}+ [sic].


\subsubsection*{\code{flix}}

\code{flix} creates bitmapped images and movies,
Section~\ref{section-animation}. Input files should have extension
\filename{flx}, and must contain a valid \code{epix} header as well as
additional code. Jay Belanger's \code{emacs} mode facilitates creation
of \code{flix} files.


\section{The Drawing Model}
\label{section-overview}

To draw a figure manually, you select a sheet of paper of appropriate
size and add paths, markers, and labels. These scene elements have
attributes (line color, line width, fill color, font size, etc.)\
affecting their appearance. 

\index{canvas}%
\ePiX's drawing model behaves similarly.  A \code{picture} command
sets the \emph{canvas} (or logical drawing area, a Cartesian
rectangle) and the true size of the final figure. A \code{begin}
command initializes the ``virtual paper'' for drawing. Subsequent
\emph{drawing commands} add objects to the scene: lines, curves,
function plots, labels, and the like. The printed appearance of scene
elements is determined by the current ``attribute state'' and
controlled by \emph{style declaration} which remain in effect until
superseded. Finally, an \code{end} command closes the figure and
prints it to the output file.

In this manual, command descriptions follow a few
conventions. Optional arguments are enclosed in square brackets. A
\code{len} argument is either a number (representing a length in
\code{pt}) or a double-quoted string containing a number and a
two-letter \LaTeX\ length unit, such as \code{"1.5pt"} or
\code{"6cm"}.  A \code{color} argument is a named primary
(\code{Red()}, \code{Cyan()}, \code{White()}, etc.), or a \code{Color}
specified by densities (\code{RGB(r,g,b)}, \code{CMY(c,m,y)}, etc.)

\begin{itemize}

\item Fill color: \code{fill(color)}, \code{nofill()}.

\item Line color and width: \code{plain([color])},
  \code{bold([color])}.

\item Text size: \code{font\_size([size])}, returns to
  \code{normalsize} if no argument is given.

\end{itemize}

\noindent A complete list of style commands is found in
Section~\ref{section-quickref}, starting on
page~\pageref{section-quickref}.

To handle three-dimensional scenes, a \code{camera} performs point
projection from a selected spatial location to the canvas. As a user,
you'll need to control relatively few of the \code{camera}'s
parameters. Keep in mind, however, that elements are added to a scene
in the same order their commands appear in the output file, and that
later elements generally cover earlier ones. Some three-dimensional
scenes require manual ordering of the input file; such ordering is
dependent on the \code{camera}'s location.

By default, \ePiX\ draws thin, solid, black lines, performs no filling
of regions, prints text in a 12~pt Roman font, and looks down the
\teximath{z}~axis from a large distance, giving orthogonal projection
on the \teximath{(x,y)}~plane. When the \code{camera} is moved, the
\teximath{z}~axis points vertically up on the page.


\section{Tutorial}
\label{section-tutorial}

This section presents sample input files side-by-side with their
output so you can compare what you write with what you'll see on the
screen or page.  \ePiX\ provides standard drawing capabilities, but
like all software has its own idioms and personality. The basic
syntax, which comes from \code{C}/\CXX, should be mostly
self-explanatory. One-line comments begin with the
string~``\verb+//+''.


%%\index{Emacs@\code{emacs}}
\index{Graphical interface}
To use the sample files interactively, you'll need working software:
\filename{bash}, \filename{emacs}, \ePiX, \filename{g++},
\filename{gv}, and a running X~server. To complete your ``GUI'',
start~\filename{gv} and select ``Watch file'' from the ``State''
menu. The loaded file will update automatically when it changes.

Create a ``scratch'' directory, \code{cd} into it, and run the command
\begin{verbatim}
  tar -zxf /usr/local/share/doc/epix/sample_src.tar.gz
\end{verbatim}
(Change the path as appropriate for your installation.)  This unpacks
copies of the sample files into your scratch directory, including all
the files mentioned in this manual. The \filename{README} file serves
as a table of contents.

Open a sample file in \filename{emacs}, compile it from the drop-down
menu (or with the keyboard command), then open the \filename{EPS} file
in~\filename{gv}. Now you're ready to follow the tutorial
interactively. A few suggested exercises are included with each
file. Naturally, as you study more files, you'll be able to make more
interesting changes on your own.


\subsection*{Basic Drawing}

The first sample, \filename{hello.xp}, contains code needed to specify
the figure's size, followed by the classic greeting. The \code{border}
command draws a box around the figure in the specified color and
width, and serves here merely to delimit the output from the
surrounding page.
\begin{verbatim}
/* -*-ePiX-*- */
#include "epix.h"      // These lines are analogous
using namespace ePiX;  // to a usepackage command.

int main()
{
  picture(P(-1,-1), P(1,1), "2 x 1in"); // corners, true size

  begin(); // ---- Figure body starts here ----

  border(Black(), "1pt"); // color, line width

  font_size("Huge");  // May be any font size, e.g. "scriptsize"
  label(P(0,0), "Hello, world!");

  end();   // ---- End figure; write output file ----
}
\end{verbatim}

\begin{center}
  \input{hello.eepic}
\end{center}

\begin{itemize}

\item Change the color and width of the
  border. (\code{RGB(r,g,b);} creates an RGB color, CMY colors are
  analogous. Named primaries are available. The densities should be
  between \teximath{0}~and \teximath{1} for ``expected'' behavior.)

\item Add \code{backing(Cyan());} after the \code{border} command. 

\item Put the command \code{crop\_ellipse();} before the \code{border}
  command. Permute the \code{crop\_ellipse} command with the
  \code{border} and \code{backing} lines, and note how the attribute
  (\code{crop}) affects objects (\code{border}, \code{backing}).

\end{itemize}


\subsubsection*{Geometric Objects}

Our next file uses simple objects to draw a 2-D house-and-sun scene.

\begin{verbatim}
int main()
{
  picture(P(0,0), P(5,2), "3.75 x 1.5in");

  begin();

  triangle(P(0.9, 1), P(3.1, 1), P(2, 1.5)); // vertices
  circle(P(4,1.5), 0.25); // center and radius

  rect(P(1,0), P(3,1)); // opposite corners

  fill(Black(0.1));     // light gray
  rect(P(2.3,0), P(2.7,0.8)); // the door
  nofill();

  bold();               // draw thicker lines
  grid(P(1.4, 0.2), P(2, 0.6), 2, 2); // corners, number of squares

  line(P(xmin(), 0), P(xmax(), 0));   // endpoints

  end();
}
\end{verbatim}

\begin{center}
  \input{objects.eepic}
\end{center}

\begin{itemize}

\item Add declarations such as \code{bold(RGB(1,0.9,0.5))} or
  \code{fill(Yellow())} to color the scene. (Color and filling are
  orthogonal attributes.)
\end{itemize}

\noindent The sample file \filename{house.flx} uses loops to draw
gradient fills of the lawn, sky, and sun, and animates a sunset.


\subsubsection*{Function Plotting}

Plotted functions must be defined in the ``preamble'', before
\code{main}. ``High-level'' elements---coordinate axes and grids, axis
labels, and graphs---are drawn with mnemonically-named commands.
\begin{verbatim}
// double = double-precision floating point
double f(double x) { return 0.75*Sin(x) - 0.25*Sin(2*x); }

int main()
{
  picture(P(-2,-1), P(2,1), "5 x 1in"); // [-2,2] x [-1,1]

  begin();
  revolutions(); // set angle units, [0,1] = one turn

  h_axis(16); // axes w/default endpts
  v_axis(4);
  h_axis_labels(4, P(0,-4), b); // shift down 4pt, align below

  plain(Red());
  plot(f, xmin(), xmax(), 120); // use 120 intervals

  end();
}
\end{verbatim}

\begin{center}
  \input{plotting.eepic}
\end{center}

\begin{itemize}

\item Define and plot some different functions; adjust the bounding
  box as necessary. (Use repeated multiplication for polynomials.
  The \code{polarplot} command graphs \teximath{r=f(\theta)}.)

\item Change the arguments to \code{h\_axis\_labels}. The first
  specifies the number of intervals to label; the second gives the
  label offset in~\code{pt}. The last puts each label below~(\code{b})
  its Cartesian location.

\end{itemize}


\subsubsection*{Multivariable Plotting}

Functions of two or three variables are defined just like functions of
one variable, but for plotting the return type must be a
point~(\code{P}), not a real number~(\code{double}). The \code{domain}
class specifies the set of inputs to plot.

\begin{verbatim}
P f(double r, double th)
{
  return P(r*Cos(th), r*Sin(th), pow(r, 3)*Cos(3*th));
}

int main()
{
  picture(P(-1,-1), P(1, 1), "2 x 1in");
  begin();
  pst_format();     // use PSTricks macros for output

  // corners ([0,1] x [0,2\pi]) and fineness
  // (8x40 rectangles, plotted at 16x120 resolution)
  domain R(P(0,0), P(1, 2*M_PI), mesh(8,40), mesh(16,120));

  camera.at(P(3,1,2)); // set the viewpoint
  arrow(P(0,0,0), P(1.25,0,0)); // coordinate axes
  arrow(P(0,0,0), P(0,1.25,0));

  plain(Blue(1.2)); // line color and width
  fill(Yellow());   // shading color
  surface(f, R);

  end();
}
\end{verbatim}

\begin{center}
  \input{plotting2.eepic}
\end{center}


\subsubsection*{Loops and Control Structures}
\label{taylor_poly}

A function can be defined by an arbitrary algorithm, and a
\code{domain} may be used to plot a family of functions for several
values of one variable.

\begin{footnotesize}
\begin{verbatim}
P sin_n(double x, double n)     // Taylor polynomial of sin x
{
  const int N((int) floor(n));  // convert n to an index bound
  const double sqx(-pow(x, 2)); // -x^2
  double val(x), summand(x);

  for (int i=1; i <= 2*N+1; i += 2)
    {
      summand *= (sqx/((i+1)*(i+2)));   // (-1)^i x^{2i+1}/(2i+1)!
      val += summand;
    }

  return P(x, val);    // return (x, y)
}

int main()
{
  picture(P(0, -1), P(6*M_PI, 1), "5 x 1in");

  begin();

  set_crop();
  bold(Green());
  plot(Sin, 0, xmax(), 120);

  domain R(P(0,1), P(6*M_PI, 41), mesh(60, 40), mesh(120, 40));
  for (int i=3; 0 <= i; --i)  // print in descending degree
    {
      bold(RGB(0.25*i, 0, 1-0.25*i)); // degree-dependent color
      plot(sin_n, R.slice2(5*i+1));   // plot for n = 5i+1
    }
  end();
}
\end{verbatim}
\end{footnotesize}

\begin{center}
  \input{plotting3.eepic}
\end{center}


\subsubsection*{Page Layout}

Page layout can be composed from sub-pages with \code{screen} objects.
In the loop body below, objects are added to the ``active''
\code{screen}, then \code{inset} into the ``canvas'', the screen
representing the entire figure.

\begin{verbatim}
P f(double u, double v)
{
  return P((u-v)*(u+v), 2*u*v, u);
}

int main()
{
  picture(P(0,0), P(2,3), "5x7.5in"); // overall size

  begin();

  domain R(P(-1,-1), P(1,1), mesh(12,12), mesh(24,24));

  for (int i=0; i<2; ++i)
    for (int j=0; j<3; ++j)
      {
        screen my(P(-3,-3), P(3,3));
        activate(my);
        border(Red(0.6), "1pt");
        // frame-dependent viewpoint
        camera.at(sph(10, (2*j+i+1)*M_PI/8, M_PI/6));

        plot(f, R);
        arrow(P(0,0,0), 2*E_1);
        arrow(P(0,0,0), 2*E_2);
        // SW corner at (i, 2-j), padded by 0.05 on all sides
        inset(P(i+0.05,2.05-j), P(i+0.95,2.95-j));
      }

  end();
}
\end{verbatim}

\begin{center}
  \input{layout.eepic}
\end{center}


\section{\CXX\ Basics}

\index{Cplusplus@\CXX|(}
An \ePiX\ source file is a \CXX~program. If you've successfully
modified and compiled the sample files, you know enough~\CXX\ to use
\ePiX. In the author's experience, \code{C}~grammar suffices for most
applications.  An excellent introduction to definitions of functions
and variables, control statements, and overall program structure is
Kernighan~and Ritchie's \emph{The \code{C} Programming Language},
second edition~\cite{KnR}.


\subsection{File Format}

\index{Emacs@\code{emacs}}
Jay Belanger's \code{emacs} mode for \code{ePiX} inserts a file
template when an empty buffer is opened with the extension~\code{xp}.
This section explains the purposes served by the template. A few
additional remarks may help you avoid basic syntax pitfalls.

A \CXX\ file consists of ``statements'', analogous to ordinary
sentences. Common types include \emph{declarations} (which
``register'' a function, variable, or type name with the compiler),
\emph{definitions} (which assign meaning to declared names), and
\emph{function calls} (which cause a named function to execute).  Most
statements in an \ePiX\ input file are function calls (``commands'').
Plain declarations are relatively rare in user files, since a
definition serves to declare any new names that it contains.

Every statement ends with a semicolon, and conventionally a file
contains at most one statement per line.  The compiler ignores
nearly all whitespace (spaces, tabs, and newlines), which should be
used freely to make files easy to read.  Other punctuation
(periods, commas, (semi)colons, parentheses, braces, and quotes)
dictates file parsing, and must adhere stringently to grammar.

An \ePiX\ file always begins with the lines
\findex{include}
\begin{verbatim}
#include "epix.h" // N.B. pre-processor directive, no semicolon
using namespace ePiX;
\end{verbatim}
The first line is analogous to a \LaTeX\ \code{usepackage} command: It
loads the contents of the ``header'' file \filename{epix.h}, importing
the names of commands provided by~\ePiX.  To avoid name conflicts,
\ePiX's commands are enclosed in a ``namespace''. For example, the
\code{label} command is actually known to the compiler as
\code{ePiX::label}. The second line above tells the compiler to apply
the prefix tacitly.


\subsection{Variables and Functions}
%% \index{Variable}\index{Variable!names}
%% \index{Function}\index{Function!names} 

Definitions of variables and functions play the same role in a figure
that macro definitions play in a \LaTeX\ document: gathering and
organizing information on which the figure depends. A variable is
defined by supplying its type, name, and initial value. By far the
most common data types in~\ePiX\ are \code{double} (double-precision
floating point number), \code{P}, and \code{int}. The name of a
variable may only consist of letters (including the underscore
character) and digits, and must begin with a letter:
\begin{verbatim}
my_var, var2,    MY_var, aLongVariableName;  // valid
my-var, 2var, \v@riable, $x, ${MY_VARIABLE}; // not valid
\end{verbatim}
Variable names are case-sensitive, and numerous (non-universal)
conventions govern the significance of capitalization. Generally, make
names descriptive but not unwieldy, and avoid language keywords (such
as \code{const}, \code{true}, \code{double}, \code{class}, or
\code{public}) and names that begin with an underscore.

A function accepts ``arguments'' and ``returns a value''. To define a
function in~\CXX, specify the return type, the name of the function,
the types of the arguments, and the algorithm by which the value is
computed from the inputs.  The code block
\begin{verbatim}
double f(double x) 
{ 
  return sqrt(1-x*x); 
}
\end{verbatim}
defines the \code{double}-valued function~\teximath{f} of one \code{double}
variable defined by the formula \teximath{f(x)=\sqrt{1-x^2}}.


\subsection{Comments}

\CXX\ has two types of comments. \code{C}-style comments, which may
span several lines, are delimited by the strings \code{/*}~and
\code{*/}. One-line comments, analogous to the \LaTeX~\code{\%}, are
begun with \code{//}. A one-line comment may appear within a multi-line
comment, but a \code{C}-style comment may not; the compiler will
mistake the first~\code{*/} it encounters as the end of the current
multi-line comment.


\subsection{Program Execution}

All the ``action'' in a \CXX\ program occurs inside the special
function~\code{main}. Running a compiled \CXX~program is viewed by the
operating system as calling the program's \code{main} function. The
return value (an \code{int}) is the program's exit status.  The
contents of the output file start with \code{begin()} and terminate
with \code{end()}.  Intervening statements constitute the \emph{body}
of the file.

In \CXX, a function may not be defined inside another function. 
Variables may be defined inside \code{main}, but functions cannot be.


\subsection{Strings and Raw Output} 

\index{string}%
\index{Backslash character}%
In \CXX, a \code{string} is a sequence of characters. Most
\code{string}s in \ePiX\ input files are \emph{literals},
double-quoted \code{string}s whose value is read from the input. In a
string literal, backslash is an escape character; a single backslash
is produced by a double backslash in the input file.  Certain letters
have special meanings when backslash-escaped, including ``\verb+\n+''
(newline) and ``\verb+\t+'' (\Key{TAB}).  Unlike~\LaTeX, \CXX~does not
require a space to separate an escape sequence from following text;
the string ``\verb+\\textwidth+'' literally represents a \LaTeX\
command, while ``\verb+\textwidth+'' is read ``\code{TABextwidth}'' by
the compiler.

\index{Self-contained figure}%
\index{Writing directly to output file}%
\index{Output file!writing directly to}%
Though not commonly needed, raw text can be printed to the output
file. The functions \code{write}, \code{pre\_write}, and
\code{post\_write} accept \code{string} arguments. \code{write} prints
its argument where the call appears in the input file. The other
functions print their arguments before or after the completed picture,
respectively. These commands must be in the file body.  As an
application, a complete \LaTeX\ \code{figure} environment (with
caption and label) can be produced by an \ePiX\ file,
Figure~\ref{fig-captioned}.

\findex{pre\_write}
\findex{post\_write}
\findex{write}
\begin{figure}
\begin{footnotesize}
\begin{verbatim}
#include "epix.h"
using namespace ePiX;

int main()
{
  picture(P(-1,-1), P(1,1), "10cm x 3in);
  begin();
  pre_write("\\begin{figure}[hbt]"); // comes before the picture

  post_write("\\caption{A \\LaTeX\\ figure.}"); // and after
  post_write("\\end{figure}");

  < ... other ePiX commands ... >
  write("%% A comment near the end, but inside the picture.");
  end();
} // End of main()
\end{verbatim}
\end{footnotesize}
\caption{Generating a self-contained figure in \ePiX.}
\label{fig-captioned}
\end{figure}


\subsection{Conditionals and Loops}

An algorithm's behavior usually depends on internal state. A
\emph{conditional statement} causes blocks of code to be executed
according to criteria. A \emph{loop} repeatedly executes a code
block, usually changing the values of variables in a predictable way,
so that the loop exits after finitely many traversals.

\texicomment
Figure~\ref{fig-gcd} illustrates conditionals and loops with
Euclid's algorithm for the greatest common divisor. Three pieces of
notation require explanation: \code{j\%i} means~``\teximath{j\pmod i}'',
\code{||} is logical~``or'', and \code{==} is ``test for equality''.
(A single~``\code{=}'' is the assignment operator.)

\index{Function!Euclidean algorithm}
\begin{figure}[hbt]
\begin{footnotesize}
\begin{verbatim}
int gcd (int i, int j)
{
  int temp(i);  // initialization syntax
    if (i==0 || j==0)
        return i+j;    // define gcd(k,0) = k

    else {
        if (j < i)     // swap them
            {
                temp = j;
                j = i;
                i = temp;
            }
        // the work is done here...
        while (0 != (temp = j%i)) // assign temp, test for zero
            {
                j = i;
                i = temp;
            }
        return i;
    }
}
\end{verbatim}
\end{footnotesize}
\caption{Euclid's division algorithm in~\CXX.}
\label{fig-gcd}
\end{figure}
\texicommentend

%% \noindent For details on conditionals, or for more advanced techniques
%% of \CXX\ programming, please consult a textbook or online tutorial.
\index{Cplusplus@\CXX|)}

%% Force previous two figures onto the same page
%%\clearpage

\section{Animation}
\label{section-animation}

\index{Animation|(} 
\ePiX\ is well-suited to the creation of mathematically accurate
animations: If a figure depends suitably upon a ``time'' parameter,
then a loop can be used to draw the entire figure for multiple time
values, yielding successive ``snapshots'' of the figure as time
progresses. The shell script~\code{flix} automates the process of
compiling a suitable input file into a collection of \filename{png}s
and assembling these frames into a \filename{mng}~or \filename{gif}
animation. GraphicsMagick is the image-handling engine.

A \code{flix} file is an \code{epix} file with two restrictions:
\begin{itemize}
  \item The \code{double} variable \code{tix()} is used as
  ``clock''.

  \item \code{main} accepts two command line arguments and sets
  \code{tix()} accordingly.
\end{itemize}
\index{Emacs@\code{emacs}}
Jay Belanger's \code{emacs} mode recognizes the file
extension~\filename{.flx} and inserts template code if an empty buffer
is opened. Creation of \code{flix} files is as easy as creation of
\code{epix} files.  The \filename{samples} directory contains a
handful of \code{flix} files that may be consulted for ideas.

By default, \code{flix} creates movies with 24~frames, in which
\code{tix()} runs from~0 to~1, and animates at
\teximath{0.08}~sec/frame.  These and other parameters can be changed
with command-line options.

A ``typical'' \code{.flx} file may take 30~seconds to a few minutes to
compile, depending on the number of frames and the complexity of each
frame. To facilitate debugging, \code{elaps} can be run on a
\code{flix} file. \code{elaps} runs in a fraction of the time, and if
\code{elaps} can't produce a viewable image, \code{flix} will surely
fail.
\index{Animation|)}


\section{Layout Tricks}
\label{section-layout}
\index{Layout}

\subsection{Stereograms}
\index{Stereograms}
Stereograms are created by drawing a single scene twice from slightly
different points of view. To create a stereogram with \ePiX, write a
``scene'' function containing the necessary drawing commands, then set
the camera appropriately and call the scene function twice, once for
each frame. 

\subsubsection*{Crossed-Eyes Stereograms}

The sample files \filename{lorenz.xp}~and \filename{twisted\_cubic.xp}
use page layout to draw crossed-eyes stereograms. For each frame,
create and activate a \code{screen}, set the camera, and call the
scene function. The frames are inset side by side in the final
picture. For a crossed-eyes stereogram, the frame from the
``rightmost'' camera position is placed on the left.

\subsubsection*{Bi-Color Stereograms}

The sample files \filename{cube.flx}~and \filename{mirrorball.flx}
contain animated red-cyan stereograms, suitable for viewing with 3-D
glasses. The technique works best with black-and-white line drawings,
but the basic approach is the same as for crossed-eyes stereograms:
Create a scene function and call it twice, setting the camera
appropriately. Layout is unnecessary since the frames are superposed.
Pen colors should be close to red and cyan, but ``optimal'' choices
depend on one's eyes and 3-D glasses.

\subsubsection*{Stereographic Movies}

Either form of stereogram can be animated in a \code{.flx} file, but
keeping one's eyes properly crossed requires practice while most
people can instantly see depth with 3-D glasses. As always when using
\code{flix}, debug the scene with \code{elaps} before compiling a
movie. If a stereoscopic effect is difficult to see in a still image,
it will be even harder to perceive in a movie.


\subsection{Inset Images}

\ePiX\ can place external images into a figure, similarly to
Rolf Niepraschk's \code{overpic} package. You provide the Cartesian
center point, the file name, and optionally the true height and/or
width of the image:
\findex{label}
\begin{verbatim}
  label(P(a,b), "\\includegraphics[width=w,height=h]{file}");
\end{verbatim}
This effect requires the \filename{graphicx} package. When compiling a
stand-alone graphic containing an external image, you must specify the
\filename{graphicx} package on the command line: \code{elaps -p
graphicx <file>}

\texicomment
\begin{figure}[hbt]
  \begin{center}
    \input{ePiX.eepic}
  \end{center}
  \caption{Visual guide to \ePiX's shell scripts}
  \label{fig-programs}
\end{figure}
\texicommentend


\texinfochapter{3}{Reference Manual}
\label{chapter-ref-man}

This chapter details \ePiX's capabilities, discussing attributes and
data types, and listing available commands in the form
\begin{verbatim}
  label(P posn, [P offset], string text, [align]);
  circle(P ctr, double rad=1, normal=P(0,0,1));
\end{verbatim}
As in Chapter~\ref{chapter-started}, function arguments are given by
type (\code{P}, \code{double}, \code{string}, etc.)\ and name, or by
name alone if the type is clear. Optional arguments are enclosed in
square brackets. A few argument types, such as \code{[align]} above,
admit only a small number of values; these cases are explained when
the corresponding command is first introduced.

A name followed by an equals sign and a value indicates a default
argument; if omitted in an input file, the compiler substitutes the
default value. Only trailing arguments may be specified this way.  For
example, if the \code{rad} argument of the \code{circle} command is
omitted, the \code{normal} argument \emph{must} be omitted as well.

In an input file, only an argument's value is given, not the type:
\begin{verbatim}
  label(P(0,0), P(0,-4), "Hello world", b);
  circle(P(0,-0.25), sqrt(2)); // use default normal
\end{verbatim}

\index{Length}%
\index{Color}%
Generally, \code{len} signifies either a number (representing a length
in \code{pt}) or a double-quoted string containing a number and a
two-letter \LaTeX\ length unit, such as \code{"1.5pt"} or
\code{"6cm"}. \code{color} represents a \code{Color} object
constructed from a named primary (\code{Blue()}, \code{Magenta(0.7)},
etc.), a color specification (\code{RGB(r,g,b)}, \code{CMY(c,m,y)},
etc.), or an operator applied to an existing \code{Color}.


\section{File Structure}

An \ePiX\ input file constitutes a short \CXX\ program. When this
program is compiled and run by one of the shell scripts, it creates a
figure file suitable for inclusion in \LaTeX.

\index{Preamble}
Like a \LaTeX\ document, an \ePiX\ file contains a \emph{preamble},
which sets up a drawing environment, and a \emph{body}, which contains
actual figure-generating commands.  The minimal file has the form
\findex{picture}
\findex{begin}
\findex{end}
\begin{verbatim}
#include "epix.h"  // N.B. no semicolon
using namespace ePiX;

int main() {
  picture(P(a1,b1), P(a2,b2), "n1 [unit1] x n2 unit2");
  begin(); // end of preamble, start of body
  end();   // end of body
}
\end{verbatim}
\ePiX\ commands are of four general types: drawing, attribute setting,
definitions (of data and functions), and operations on existing
objects. Except as noted below, drawing and attribute commands must
appear in the body, between \code{begin()}~and \code{end()}. Function
definitions must come in the preamble, before \code{main()}. Data
definitions may appear in the preamble or body.


\subsection*{Output Format}
\index{Output format}

The \code{end()} command writes the output file to \code{stdout},
using \code{eepic} macros by default. The shell scripts redirect
\code{stdout} to an appropriate disk file.

The attribute-setting command \code{pst\_format()} causes the output
file to be written using \code{PSTricks} macros when \code{end} is
called. Similarly, \code{tikz\_format()} causes the file to be written
using \code{tikz} macros, and \code{eepic\_format()} causes the file
to be written using \code{eepic} macros. These commands may appear
anywhere in the figure body. There is no reason for a file to contain
more than one such command.

The output format may be selected on the command line, overriding any
explicit request in the input file. Supplying \code{epix},
\code{elaps}, or \code{flix} with one of the options \code{--pst},
\code{--tikz}, or \code{--eepic} is tantamount to issuing a
\code{pst\_format()} (etc.)\ command just before the end of the
file. These options are listed in decreasing precedence. If more than
one is given, the ``strongest'' applies, regardless of the command
line order.

A figure may be written directly to a specified disk file, in a format
unaffected by the command line flags above. The command
\begin{verbatim}
  print_eepic("file.tex");
\end{verbatim}
writes the figure to the named file using \code{eepic}
macros. Analogous commands exist for the formats \code{pst}~and
\code{tikz}. These commands perform an immediate action. Consequently,
the command's location in the input file is significant, and multiple
commands may appear in a single file, so long as distinct file name
arguments are provided. Applications include writing the same figure
in multiple formats, or creating successive ``snapshots'' of a lengthy
computation.


\section{Picture Size and Aspect Ratio}

\index{Picture!size}%
\index{Picture!offset}%
\index{Aspect ratio}%
The \code{picture} command specifies a figure's logical and true
sizes. In the ``minimal file'' snippet above, the \emph{canvas} is the
Cartesian rectangle \teximath{[a_1,a_2]\times[b_1,b_2]} whose corners
are given. Either pair of opposite corners is acceptable, but
confusion is less likely when the SW~and NE corners are given, in this
order.

\index{Length}%
The true (printed) width and height are read from \code{picture}'s
third argument, a ``size \code{string}'' such as \verb+"4cm x 1in"+,
containing a number and optional \LaTeX\ length unit, an~\code{x},
another number, and a mandatory length unit. Spaces may be used for
legibility.  The mandatory length unit becomes the \LaTeX\
\code{unitlength}.  The argument \verb+"4cm x 1in"+ creates a picture
4~cm wide and 1~in high, but internally converts 4~cm to inches.  
Recognized length units are \code{pt} (points, the default), \code{cm}
(centimeters), \code{in} (inches), \code{mm} (millimeters), \code{pc}
(picas), and \code{bp} (big points).  
(\code{1in = 2.54cm = 72bp = 72.27pt}, \code{1pc = 12pt}.)

The logical and true sizes may be defined separately with the commands
\findex{bounding\_box}
\findex{picture}
\findex{unitlength}
\begin{verbatim}
  bounding_box(P(a1,b1), P(a2,b2));
  picture(n1, n2);
  unitlength(len); // len a length string, not a double
\end{verbatim}
One way or another, the logical and true dimensions \emph{must} have
positive values when the \code{begin()} command is issued.

\LaTeX\ treats the contents of a picture environment as a single box,
aligned by default on its lower left corner. An \code{offset} command
accepts two \code{double} arguments or a size \code{string} and shifts
the page location accordingly. The command
\code{offset("0.25 x -0.5cm")} shifts the picture right
\teximath{0.25}~cm and down \teximath{0.5}~cm.  If the
\code{unitlength} is 1~cm, the command \code{offset(0.25, -0.5)} has
the same effect.

A non-zero \code{offset} causes a picture's contents to appear in a
location where \LaTeX\ does not expect them. This can be useful in a
\LaTeX\ document, but should be avoided when compiling a stand-alone
image, since \code{dvips} may crop according to rules of its own.


\index{Picture!aspect ratio}
\index{Picture!bounding box}
The canvas's corners are \code{(xmin(),ymin())} and
\code{(xmax(),ymax())}, while its width and height are
\code{xsize()}~and \code{ysize()}. The canvas is a virtual, advisory
data structure; its dimensions are not directly related to the
figure's printed size, and picture elements may lie outside the
canvas. Affine scaling maps the canvas to the page when the output
file is written.
\begin{center}
\input{scaling.eepic}
\end{center}
\label{page-scaling}
The figure's aspect ratio is controlled by sizing the canvas.  The
aspect ratio is ``true'' if the canvas and page rectangles are
geometrically similar, e.g., if both boxes are 1.5~times as wide as
they are tall.


\section{Color}
\label{section-color}

\index{Color|(} 
\ePiX\ provides a \code{Color} data type. Four models are implemented:
\code{RGB}, \code{CMY}, \code{CMYK}, and \code{Gray}.  A \code{Color}
holds primary color ``channels'', each carrying an intensity between
0~(no color) and~1 (full saturation).  ``No color'' means black in
\code{RGB} and \code{Gray}, white in \code{CMY(K)}.  \code{RGB},
\code{CMY}, and~\code{CMYK} colors are written to the output file as
commands in the corresponding model. \code{Gray} shades are written
in~\code{RGB}.

Operations on red-green-blue colors are described below. Functionally,
colors are converted to~\code{RGB}, operated upon, then converted back
to the original model. Conversions are as described in Uwe Kern's
\filename{xcolor} manual~\cite{xcolor}.


\subsection{Constructors}

Each color model has a ``constructor'' creating a color of specified
densities. The \code{RGB}~and \code{CMY} models have named ``primary''
constructors; the density argument is optional and defaults to~1.
\findex{RGB}
\findex{Red}
\findex{Green}
\findex{Blue}
\findex{White}
\findex{Black}
\index{Color!primitive}
\begin{verbatim}
  // red-green-blue colors
  RGB(double r=0, double g=0, double b=0);
  Red(d=1);   Green(d=1);   Blue(d=1);
  White(d=1);   Black(d=1);
\end{verbatim}
\findex{CMY}
\findex{Cyan}
\findex{Magenta}
\findex{Yellow}
\findex{CMY\_White}
\findex{CMY\_Black}
\findex{CMYK}
\findex{Gray}
\begin{verbatim}
  // cyan-magenta-yellow colors
  CMY(double c=0, double m=0, double y=0);
  Cyan(d=1);   Magenta(d=1);   Yellow(d=1);
  CMY_White(d=1);   CMY_Black(d=1);

  // gray
  Gray(double d=0); // equivalent to RGB(d, d, d)
\end{verbatim}
Each \code{CMY} constructor has a corresponding \code{CMYK} function,
e.g.\ \code{CyanK()}~or \code{CMYK\_White()}.

\index{Color!density|(}%
Though color densities lie between 0~and 1, \ePiX's primary color
constructors take arguments mod~4, viewed as elements
of~\teximath{[-2,2]}. Consider \code{Red(d)}, ``red with
density~\teximath{d}''. For integer values of~\teximath{d}, the
constructor has the following meanings: \code{Red(0)}~is black,
\code{Red(1)}~is red, \code{Red(2)}=\code{Red(-2)}~is white, and
\code{Red(-1)}~is anti-red, or cyan. For non-integer~\teximath{d}, the
constructor interpolates between the bracketing integer values,
Figure~\ref{fig-densities}. Other primary constructors work
analogously.

\begin{figure}[hbt]
  \begin{center}
    \input{color_density.eepic}
  \end{center}
  \caption{The red primary constructor.}
\label{fig-densities}
\end{figure}

Let \teximath{\delta:[-2,2]\to[0,1]} be the piecewise-linear function
that converts real numbers to primary color densities,
\teximath{\clip:\mathbf{R}\to[0,1]} the clipping function.  The
non-primary constructor \code{RGB(r, g, b)} computes the density of
the red channel as
\teximath{\clip\bigl(\delta(r)+\delta(-g)+\delta(-b)\bigr)}; the green
and blue channel densities are computed similarly. Every color can be
created with arguments between 0~and 1, but the constructor accepts
arbitrary real arguments and returns colors varying ``continuously and
periodically''.
\index{Color!density|)}%

\index{Color!transparency} 
Colors possess an ``alpha-channel'' for transparency. Except as
supported by the output format, this feature is a stub.


\subsection{Color Operations}

\index{Color!blencing}%
Colors can be scaled, inverted, blended, superposed, and filtered. In
the code below, \code{tint}~is a \code{Color}. The effect of each
operation is described in the \code{rgb} model; the visual result is
the same in all models.
\findex{Color.invert}
\findex{Color.blend}
\findex{Color.superpose}
\findex{Color.alpha}
\findex{Color.filter}
\begin{verbatim}
  tint *= double c; // multiply channels by c, rebuild
  tint.invert();    // (r, g, b) -> (1-r, 1-g, 1-b)
  tint.blend(Color col, double t);  // (1-t)*tint + t*col
  tint.superpose(col); // add channels, then clip to [0,1]
  tint.alpha(double d); // set alpha channel to clip(d)

  tint.filter(col); // return min density in each channel
\end{verbatim}
Except for \code{filter}, these operators modify their object. 
\code{filter} simulates the effect of viewing \code{col}
through a transparent sheet of \code{tint}, and returns a new
\code{Color} object having the same model as \code{tint} without
modifying \code{tint} itself.


\section{Scene Attributes}

Objects in a scene sit in 3-dimensional space. A \emph{camera}
maps objects to the \emph{active screen}. The active \code{screen}
may, in turn crop its contents. Each drawing command creates an
object, ``photographs'' it, and adds the image to the active
\code{screen}. This section describes the \code{Camera} and
\code{screen} classes and a few associated concepts.


\subsection{Angular Mode}
\index{Angle!units}

By default, angles are measured in \code{radian}s. Two other angular
modes are available: \code{degrees}~and \code{revolutions}. The
angular mode is set with a named command, \code{radians()},
\code{degrees()}, or \code{revolutions()}. The current angular mode
affects all trigonometric functions and operations.


\subsection{The Camera}
\index{Camera|(emph}
\index{Screen|(}

Three-dimensional scenes are drawn on flat paper by applying a
mathematical transformation. By default, \ePiX\ uses \emph{point
projection}, the technique used by art students when they trace on a
window with grease pencil, Figure~\ref{fig-camera}.

\begin{figure}[hbt]
  \begin{center}
    \input{camera.eepic}
  \end{center}
  \vspace*{-0.25in}
  \caption{Point projection.}
  \label{fig-camera}
\end{figure}

\index{Camera!viewpoint}
\index{Viewpoint}
\ePiX\ depicts a Cartesian world by projecting mathematically to a
screen plane, then affinely scaling to a printed page. The camera,
which maps the world to the screen, consists of a \emph{body} (data
that determines the position and orientation of the camera), a
\emph{lens} (the actual mapping to the screen plane), and a
\emph{filter} (a color through which the scene is viewed).


\subsubsection*{Body}

The camera's spatial orientation is described by a triple of mutually
perpendicular unit vectors.  In memory of happy days at the beach,
these vectors are called \emph{sea}, \emph{sky}, and~\emph{eye}. The
screen plane is parallel to the sea-sky plane; the sea vector points
horizontally to the right, sky points vertically upward. The eye is
their cross product, which points directly at the viewer.

The sea-sky-eye basis is located at the \emph{viewpoint}, the camera's
spatial location. The \emph{target}, the origin of the screen plane,
lies on the line through the viewpoint in the direction opposite the
eye vector.  The distance from the viewpoint to the target is the
\emph{range}.  The \emph{camera sphere} is centered at the viewpoint
and passes through the target.  The orientation, viewpoint, target,
and range completely (and redundantly) determine the camera's
geometric situation in the world.


\subsubsection*{The Lens}
\index{Camera!lens}

A \emph{lens} maps the world to the screen.  \ePiX\ comes with four
lenses: \emph{shadow} (the default), \emph{orthogonal},
\emph{fisheye}, and \emph{bubble}. Each lens simulates the appearance
of world objects as seen by an observer at the viewpoint. The shadow
lens is point projection from the viewpoint to the screen plane. The
orthogonal lens projects from infinite distance.  Each of the other
lenses performs radial projection to the camera sphere, then maps the
sphere to the screen plane; the fisheye lens does orthogonal
projection (so the entire image lies inside the disk centered at the
target whose radius is the range) while the bubble lens does
stereographic projection from the target's antipode.
\index{Screen|)}


\index{Camera!manipulation|(}
An input file normally uses a single camera, so a global object named
\code{camera} is defined automatically. At the start of a figure, the
\code{camera} looks down on the \teximath{(x_1,x_2)}-plane from a
distant point on the \teximath{x_3}-axis. The resulting view,
essentially projection along the axis, is suitable for 2-dimensional
figures.

The \code{camera} is manipulated similarly to a real camera:
\findex{Camera.at}
\findex{Camera.look\_at}
\findex{Camera.range}
\findex{Camera.focus}
\findex{Camera.tilt}
\findex{Camera.pan}
\findex{Camera.roll}
\findex{Camera.clip\_range}
\begin{verbatim}
  camera.at(P posn);          // set viewpoint to posn
  camera.look_at(P targ);     // set target to targ
  camera.range(double dist);  // fix target, move viewpoint
  camera.focus(double dist);  // fix viewpoint, move target
  camera.tilt(double angle);  // rotate about the sea axis
  camera.pan(double angle);   // rotate about the sky axis
  camera.roll(double angle);  // rotate about the eye axis
  camera.clip_range(dist);    // cull scenery closer than dist
\end{verbatim}
Explicitly setting the viewpoint or target makes the sky vector
parallel to the projection of the \teximath{x_3}~axis when possible;
otherwise the \teximath{x_2}-axis is used. Setting the range or focus
moves the \code{camera} parallel along the eye line. Each command
re-sizes the image; note that increasing the focus \emph{enlarges} the
image. The three rotation operations fix the viewpoint, but only roll
fixes the target.


\subsubsection*{Filtering and Color Separation}
\index{Color!separation}
\index{Camera!color separation}

The \code{camera} has a \emph{filter} through which all scene
\code{Color}s pass. For uniformity, the filter should only be set at
the start of the file body, before any visible elements have been
placed in the scene.  The filter has two primary uses---model
conversion and color separation---but can also be used for special
effects.

There is a \code{Neutral} color, for which filtering has no effect at
all. The \code{camera}'s default filter is \code{Neutral}. In
addition, each color model has a \emph{neutral} representative,
\code{RGB\_Neutral}, etc. Passing \code{Color}s through a neutral
filter preserves their appearance but converts them to the neutral
\code{Color}'s model.

Filtering can be used to ``split'' a \code{Color} into primary
constituents. For \code{RGB} channels, the primaries themselves are
suitable filters. For \code{CMYK} channels, there are special
\emph{process} filters, named \verb+C_Process+, etc.
\begin{verbatim}
  camera.filter(CMY_Neutral());  // convert all colors to CMY
  camera.filter(Gray_Neutral()); // black and white photography
  camera.filter(Green());        // the green channel
  camera.filter(M_Process());    // the magenta channel
  camera.filter(Red(1.4));       // la vie en rose?
\end{verbatim}
\index{Camera!manipulation|)}
\index{Camera|)}


\subsection{Clipping}

\index{Clipping|(}%
Two operations handle elements lying far from the target.
\emph{Cropping} culls elements based on their screen location, and is
described later. \emph{Clipping} removes objects whose spatial
location lies outside the ``clip region''.

Initially, the clip region is a very large box centered at the
origin. Commands are provided to resize this box, keeping the faces
parallel to coordinate planes.
\findex{clip\_box}
\findex{clip\_to}
\begin{verbatim}
  clip_box(P pt1, P pt2); // opposite corners
  clip_box(P pt);         // opposite corners pt and -pt
  clip_to (P pt);         //                  pt and P(0,0,0)
  clip_box();             // very large box
\end{verbatim}
Additional ``clippers'' may be added manually. Individual clipping
planes are specified by a point and inward-pointing normal vector.
Parallel planes cutting a ``slice'' or ``slab'' are described by a
location, direction, and a distance. The location lies halfway between
the clipping planes, and the planes' normals point toward the location.
\findex{clip\_face}
\findex{clip\_slice}
\findex{clip\_restore}
\begin{verbatim}
  clip_face(P loc, P perp); // perp points inward
  clip_slice(loc, perp);    // very close parallel planes
  clip_slice(loc, perp, dist); // planes separated by dist
  clip_restore();  // remove manually-added clip planes
\end{verbatim}
Like \code{clip\_restore}, the \code{clip\_box} and \code{clip\_to}
commands remove all user-specified half spaces. 
\index{Clipping|)}


\subsection{Screens and Page Layout}
\index{Screen|(}
\index{Screen!active}
\index{Layout}

By default, drawing occurs in the Cartesian coordinate system of the
\code{canvas}. However, insets and subfigures are most naturally
composed in ``local'' coordinates, then placed into their final
location as a unit, a process called \emph{layout}. \ePiX\ implements
layout with the \code{screen} class.


\subsubsection*{Screens}
\index{Bounding box}

A \code{screen} is a Cartesian plane containing a distinguished
coordinate rectangle, its \emph{bounding box}. The \code{canvas} of a
figure is a \code{screen}, as is the internal representation of the
printed page.  At any point of a file body, some \code{screen} is
\emph{active}, or ``open for drawing''. A new \code{screen} is created
from a pair of opposite corners, but is not used until explicitly
activated. 
\findex{activate}
\findex{deactivate}
\begin{verbatim}
  screen subfig(P(a1,b1), P(a2,b2)); // [a1, a2] x [b1, b2]
  activate(subfig);   // open subfig for drawing
  // commands to draw in subfig
  deactivate(subfig); // re-activate the previous screen
\end{verbatim}

The \code{canvas} is automatically drawn at the \code{end} of the
input file. The contents of other \code{screen}s must be incorporated
into the \code{canvas} explicitly with an \code{import}~or
\code{inset} command. \code{import} affinely scales the bounding box
of the ``child'' \code{screen} to the \code{canvas}. \code{inset}
specifies corners in the ``parent'' where the child is
placed.
\findex{import}
\findex{inset}
\begin{verbatim}
  import();           // active screen to canvas
  import(subfig);     // subfig to canvas
  inset(P sw, P ne);  // active screen to specified rectangle
  inset(subfig, sw, ne); // subfig to specified rectangle
  inset(subfig);      // inset to subfig's corners
\end{verbatim}


\subsubsection*{Extracting}

The portion of a \code{screen} enclosed by a specified rectangle can
be \emph{extracted}. Each of the commands
\findex{extract}
\findex{canvas}
\begin{verbatim}
  subfig.extract(P sw, P ne);
  canvas().extract(P sw, P ne);
\end{verbatim}
returns the \code{screen} having stated corners and containing part of
\code{subfig} or the \code{canvas}, respectively. Elliptical and
diamond-shaped portions of a \code{screen} may be extracted.  The
contents occupy the ellipse or diamond inscribed in the given
rectangle.
\findex{extract\_diamond}
\findex{extract\_ellipse}
\begin{verbatim}
  subfig.extract_ellipse(sw, ne);
  canvas().extract_diamond(sw, ne);
\end{verbatim}
The \code{screen} returned by \code{extract} (or its variants) may be
used just like a manually-constructed \code{screen}, for example to
construct a magnified inset. 


\subsubsection*{True-Size Drawing}
\label{truesize}

On occasion, it is necessary to draw at known page size, independently
of the unit length. For 2-dimensional drawing in the
\teximath{(x_1,x_2)}-plane with the \code{camera} in its default
location, the coordinates (horizontal, vertical, or both) of a
\code{screen} may be interpreted as true~pt offsets relative to a
specified point rather than as Cartesian coordinates.
\findex{inlay}
\findex{inlay\_horizontal}
\findex{inlay\_vertical}
\begin{verbatim}
  inlay(subfig, P loc);
  inlay_horizontal(subfig, loc);
  inlay_vertical(subfig, loc);
\end{verbatim}
Each command maps the origin of \code{subfig} to the specified
location in the active \code{screen}. The first performs affine
scaling so that one Cartesian unit in \code{subfig} maps to one
true~pt on the page.  The second and third commands perform this
scaling only in the named direction.

\findex{pt\_to\_screen}
The \code{inlay} functions are suitable only for 2-dimensional
effects. For 3-dimensional true-size drawing, the function
\code{pt\_to\_screen} converts a true length of 1~pt into Cartesian
coordinates in the active \code{screen}. For example, the command
\code{line(P(0,0), P(pt\_to\_screen(12),0))} draws a line segment
12~pt long with its left endpoint at the Cartesian origin. True size
drawing is unusual in an input file, but has definite uses in library
code; \ePiX's right angle marker is a typical application.

Either type of true-size drawing works as described when drawing in
the \code{canvas}, and \emph{only} when drawing in a \code{screen}
imported to the \code{canvas} at ``true Cartesian'' size (one
\code{screen} unit maps to one Cartesian unit in the \code{canvas}).
In other situations, the true size coordinate(s) will be scaled by
additional affine transformation(s).


\subsubsection*{Cropping and Decorations}

\index{Cropping}
A \code{screen} has a \emph{crop mask}: rectangular (default),
elliptical, or diamond-shaped. If cropping is switched on in the
active \code{screen}, added elements are cropped to the crop mask.
For example, cropping by the default crop mask in the \code{canvas}
ensures the figure lies inside the printed region allocated by
\LaTeX.

The \code{crop} commands below affect the active \code{screen}
only. Cropping is not a ``global'' attribute, but must be set or unset
explicitly for each \code{screen}. By default, cropping is switched
off.
\findex{crop}
\findex{crop\_ellipse}
\findex{crop\_diamond}
\findex{set\_crop}
\begin{verbatim}
  set_crop();      // activate cropping
  set_crop(false); // deactivate cropping
  crop_diamond();  // <>-shaped mask, set crop state
  crop_ellipse();  // etc...
  crop();          // crop active screen's contents
\end{verbatim}
Normally the crop mask is inscribed in the \code{screen}'s bounding
box. The commands \code{crop\_diamond} and \code{crop\_ellipse} also
accept a pair of \code{P}~arguments, which are treated as corners of
the crop box. 

The active \code{screen} may be decorated with a \code{border}
(outline of specified color and line width) and \code{backing} (solid
fill color). The crop mask determines the shape of the \code{border}
and \code{backing}.
\findex{border}
\findex{backing}
\begin{verbatim}
  border(color, len);
  backing(color);
\end{verbatim}
Identically-named member functions can be applied to an arbitrary
\code{screen} whether or not it is active:
\begin{verbatim}
  scr.crop_rectangle().backing(Blue(1.8));
  canvas().crop();
\end{verbatim}


\subsubsection*{Affine Maps}
\index{Affine map}

A \emph{plane affine map} has the form \teximath{T(x) = Ax+b} for some
invertible \teximath{2\times2} matrix~\teximath{A} and a constant
vector~\teximath{b}. \ePiX's \code{affine} class allows affine maps to
be built and applied to \code{screen}s' contents.

An affine map is uniquely determined by the images of three
non-collinear points.  The constructor returns the affine map sending
the points \teximath{(1,0)}, \teximath{(0,1)}, and~\teximath{(0,0)} to
\code{pt1}, \code{pt2}, and~\code{pt0}, respectively.
\findex{affine}
\begin{verbatim}
  affine af(P pt1, P pt2, P pt0=P(0,0));
\end{verbatim}
To emphasize, the arguments are locations, not displacements, and the
image of the origin comes last.  There is also a \code{void}
constructor (taking no arguments) which returns the identity map.

To facilitate construction of affine maps, an existing \code{affine}
may be post-composed with a variety of ``elementary'' affine
transformations. In the commands below, \code{th} is an angle (in
current units), \code{sc} is a non-zero \code{double}, and \code{ctr}
is a point fixed by the composing transformation. In each command,
\code{ctr} defaults to \teximath{(0,0)}.
\findex{shift}
\findex{rotate}
\findex{reflect}
\findex{h\_scale}
\findex{v\_scale}
\findex{scale}
\findex{h\_shear}
\findex{v\_shear}
\findex{postcomp}
\begin{verbatim}
  affine af; // the identity map
  af.shift(P arg);    // translate by arg
  af.rotate(th, [ctr]);  // counterclockwise rotation about ctr
  af.reflect(th, [ctr]); // reflect across line through ctr

  af.h_scale(sc, [ctr]); // horizontal scaling
  af.v_scale(sc, [ctr]); // vertical scaling
  af.scale(sc, [ctr]);   // dilatation

  af.h_shear(sc, [ctr]); // shear preserving horizontals
  af.v_shear(sc, [ctr]); // shear preserving verticals

  af.invert();           // the inverse
  af.postcomp(affine f); // post-compose with f
  af(f);                 // pre-compose, af not modified
\end{verbatim}
A non-invertible \code{affine} can be created \emph{only} by shearing
or scaling with an extremely large or extremely small argument, or by
supplying three collinear points to the constructor.  Calling
\code{invert} on a non-invertible map merely issues a warning and
performs no action.

An \code{affine} may be applied to a \code{screen}'s current contents.
The \code{screen} class has eight member functions (\code{shift}
through \code{shear}) with syntax identical to the affine map
functions. In addition, an arbitrary \code{affine} may be applied to a
\code{screen}:
\findex{apply}
\begin{verbatim}
  scr.shift(arg); // shift scr's contents by arg; etc.
  scr.apply(f);   // apply f to scr's contents
\end{verbatim}
To apply a composition of several maps to a \code{screen}, it's best
to build an \code{affine} map by composition, then \code{apply} the
map. Composing \code{affine}s is cheap; \code{apply}ing an
\code{affine} is costly in proportion to the number of elements in the
\code{screen}.

Applying an \code{affine} to a \code{screen} has no effect on the
bounding box, \code{border}, or \code{backing}, and may move elements
outside the bounding box even if cropping is active. To ensure a
\code{screen}'s contents lie inside the bounding box, \code{crop} the
\code{screen} \emph{after} \code{apply}ing the \code{affine}(s).

The sample file \filename{inverse.xp} uses \code{affine} maps to
depict branches of inverse functions in one variable,
\filename{symmetry.xp} depicts the permutation group~\teximath{S_3}
by its action on a regular hexagon.
\index{Screen|)}


\section{Drawing Attributes}

\ePiX\ maintains drawing states for filled regions, paths, and text
objects.


\subsection{Filled Regions}

\index{Filling}
\index{Path|(}
\index{Path!filled}
Filling is either on or off. When filling is active, closed paths are
filled with the current \emph{fill color}.
\findex{fill}
\findex{nofill}
\begin{verbatim}
  fill();   // turn filling on
  fill(color); // turn filling on, specify color
  nofill(); // turn filling off, same as fill(false);
\end{verbatim}


\subsection{Paths}

\index{Drawing style|(}
\index{Line style|(}
Paths and borders of filled regions are drawn with two \emph{pen}s,
each described by color and line width. The \emph{line pen} draws
all path-like objects. The \emph{base pen} does nothing unless it is
wider then the line pen, in which case it draws an ``underlayer'' or
``border'' on the line pen.
\findex{pen}
\begin{verbatim}
  pen(len);          // set line pen width
  pen(color);        // set line color, keep width
  pen(color, len);   // set color and width
\end{verbatim}
As usual, \code{len} may be either a length string or a \code{double},
interpreted as a width in~\code{pt}.  There are \code{base} functions
with the same signatures and analogous meanings for the base pen.

The line width can be set with named declarations; the optional
argument sets the line color:
\findex{plain}
\findex{bold}
\findex{bbold}
\begin{verbatim}
  plain([color]); // 0.4pt
  bold([color]);  // 0.8pt
  bbold([color]); // 1.6pt
\end{verbatim}


\subsubsection*{Path Style}

\index{Path!style}
By default, path-like objects are drawn with solid lines. Dashed,
dotted, and free-form path style patterns are also available. (The
\code{base} underlayer is always solid.) The page length of a pattern
defaults to 12~pt, but can be set.
\findex{line\_style}
\findex{dash\_size}
\begin{verbatim}
  line_style(string);
  dash_size(len=12); 
\end{verbatim}
The path style is set with a WYSIWYG string of dashes, spaces, and
periods, representing a pattern of dashes, gaps, and dots. In the
sample styles below, the repeating units have the same page size,
12~pt. 
\begin{center}
  \input{pathstyle.eepic}
\end{center}
For brevity and uniformity, named commands are provided.
\findex{solid}
\findex{dashed}
\findex{dotted}
\begin{verbatim}
  solid(); "-"     dashed(); "-  -"    dotted(); " . "
\end{verbatim}
A sequence of \teximath{n}~dashes, spaces, and dots corresponds to a
dash/dot pattern in an interval divided into \teximath{n}~subintervals
of equal length. If the \teximath{i}th character is a dash or space,
the \teximath{i}th~subinterval is drawn solid or empty,
respectively. If the \teximath{i}th~character is a period, a dot is
placed at the midpoint of the \teximath{i}th~subinterval.

This pattern is applied to a path-like object as follows. The page
length of each edge is divided by the current \code{dash\_size} and
the ceiling (next largest integer) taken. This many copies of the
current path style are scaled onto the edge. The first and last
characters are adjacent in repeated units.

There are three inequivalent ways to adjust the dash length/dot
spacing in a non-line path: Change the \code{dash\_length}, create a
path with a different number of points, or use a longer, repetitive
pattern.  For best results, the style string should not be longer than
about a dozen characters.

If you need several dashed/dotted line styles in multiple figures,
it's best to define a custom header instead of hard-coding line
styles. See Section~\ref{section-extensions} for guidance.
\index{Drawing style|)}
\index{Line style|)}
\index{Path|)}


\subsection{Text Objects}

\index{Label|(emph}
\index{Label!alignment}
\index{Marker}
Two types of textual element may appear in a file: \emph{labels} (text
boxes) and \emph{markers} (\LaTeX\ symbols). A marker occupies a box
of zero size, and is placed at a specified Cartesian location. A label
has typographical size, and is usually offset from its Cartesian
location. For placement, an \emph{alignment point} is attached to each
label, and Cartesian coordinates position the alignment point.

\subsubsection*{Labels}

A label is printed as a \LaTeX\ box.  By default, the alignment point
is its reference point, the intersection of the left edge and the
baseline, which is used by \LaTeX\ to position the box on the
page:\input{basepoint.eepic}

\index{Label!alignment|(}
The alignment point may be \emph{offset} manually by a specified
number of \code{pt}. Additionally, a label's location with respect to
the alignment point can be chosen with an optional \LaTeX-style
alignment option. This scheme allows labels to be placed easily where
they will not overlap other parts of the figure, and ensures labels
stay properly positioned when the size or aspect ratio of a figure
changes.
\findex{label}
\begin{verbatim}
  label(P posn, P offset, string msg, [align]);
  label(posn, msg);
\end{verbatim}
The first two components of the \code{offset} argument are numbers
of~\code{pt} to shift the alignment point right and up.  The optional
\code{align} argument may be one---or an appropriate pair---of
\code{t}, \code{b}, \code{r}, or~\code{l} (top, bottom, right, left),
or~\code{c} (center). These alignment options specify the position of
the label \emph{relative to the Cartesian location}~\code{posn},
namely they work \emph{opposite} to the way they work in \LaTeX.
\begin{center}
  \input{alignment_lr.eepic}\qquad
  \input{alignment.eepic}
\end{center}
\index{Label!alignment|)}

The \code{msg} argument is usually a snippet of \LaTeX\ code enclosed
in double quotes. \CXX~treats ``\verb+\+'' as an escape character, so
a double backslash is needed in the source to get a single backslash
in the output. For example,
\begin{verbatim}
  label(P(0,0), P(2,-1), "$\\rho = \\sin\\theta$", br);
\end{verbatim}
positions an alignment point \code{2pt}~right and \code{1pt}~below the
(Cartesian) origin, and typesets the equation
\teximath{\rho=\sin\theta} below and to the right.

\index{Label!rotated}
\index{Angle!units}
Labels can be rotated; the (counterclockwise) angle is set in current
angle units with the command \code{label\_angle(theta)}. For example,
a rotation angle of~90 degrees prints labels along a vertical
axis. Though label rotation has legitimate uses, it can make labels
more difficult to read, defeating their purpose. Do not use rotated
labels merely because they are available. As a practical matter, if an
output file contains rotated labels, the enclosing document must use
the \code{rotating} package. \code{elaps} automatically loads this
package.

\index{Affine map!action on labels}
An \code{affine} behaves moderately intelligently when acting on a
marker or label. The mapping is applied to the label's location, and
the ``linear part'' is used to adjust the offset and label angle. No
attempt is made to treat alignment. For best results, if a
\code{screen} will have an \code{affine} applied to it, label
positions should be fine-tuned only with offsets, not alignment
arguments.  In all situations, the font itself is unchanged; no
attempt is made to print sheared, scaled, or reflected text. The
sample file \filename{inverse.xp} illustrates the effect of affine
maps on labels.


\subsubsection*{Fonts and Type Size}

\index{Fonts}
\index{Label!fonts in}
By default, the font in an \ePiX\ figure is that of the enclosing
document. The font size and face are changed with
``declaration-style'' commands such as
\findex{font\_size}
\findex{font\_face}
\begin{verbatim}
font_size("Large");
font_face("sc");
\end{verbatim}
The argument of \verb+font_size+ is a valid \LaTeX\ size. If no
argument is given, \code{normalsize} is understood. \verb+font_face+
accepts a two-letter string, appended to the string ``\code{text}'' to
give a \LaTeX\ font declaration command (``\code{textsc}'' above).
Finer-grained control is accomplished by placing \LaTeX\ commands
into the label text.


\subsubsection*{Label Attributes and Masked Labels}

\index{Label!masked}%
Each label command has a corresponding ``mask'' version
(\code{masklabel}) that draws an opaque rectangle under the label
text. The mask size exceeds the label's size by an amount of
\emph{padding}, and the mask is itself surrounded by a rectangular
\emph{border}, of specified color and width.
\findex{label\_color}
\findex{label\_mask}
\findex{label\_pad}
\findex{label\_border}
\findex{no\_label\_border}
\begin{verbatim}
  label_color(color);         // set label text color
  label_mask(color=White());  // set mask color
  label_pad(string len);

  label_border(color, [len]); // set color (and line width)
  label_border(len);          // set line width only
  no_label_border();          // turn off label borders
\end{verbatim}


\subsubsection*{Markers}

\index{Marker!types of|(}
\ePiX\ markers are obtained with \verb+marker(P pt, <MARKER TYPE>);+
(Table~\ref{table-marker}.)
Several ``dot-like'' marker types are available by name: \input{dots.eepic}
\findex{spot}
\findex{dot}
\findex{ddot}
\findex{box}
\findex{bbox}
\findex{ring}
\findex{circ}
\begin{verbatim}
  spot(P pt);  dot(P pt);   ddot(P pt);
     ---       box(P pt);   bbox(P pt);  
  ring(P pt);  circ(P pt);     ---
\end{verbatim}
A \code{circ} is filled with the current \code{mask\_color}, and a
\code{ring} is ``hollow''.  Each dot-like marker can also be called
with label syntax, generating a labeled marker with one command.
\begin{verbatim}
  dot(P posn, offset=P(0,0), msg="", align=none); // etc.
\end{verbatim}

\findex{dot\_size}
By default, \code{spot}~and \code{ring} are 4~pt in diameter;
\code{dot}, \code{box}, and~\code{circ} are 3~pt in diameter;
\code{ddot}~and \code{bbox} are 2~pt in diameter. At arbitrary
diameter, a \code{spot} is \teximath{4/3}~the diameter of a~\code{dot}
and a~\code{ddot} is \teximath{2/3}~the diameter.  The command
\code{dot\_size(diam=3)} sets the diameter of a \code{dot}, and hence
the size of all dot-like markers.

\begin{table}[tb]
  \begin{center}
    \input{marker.eepic}
  \end{center}
  \caption{\ePiX's \code{marker} types.}
  \label{table-marker}
\end{table}
\index{Marker!types of|)}


\subsubsection*{Reminders}

\index{Label!offset}%
When constructing and placing a label,
\begin{itemize}
\item Offsets are specified in \code{pt} (true length), not Cartesian
  units: A label's location relative to its alignment point should not
  depend on the logical or printed size of the figure.

\item The label text is enclosed in double quotes (the single
  character~\code{"}), and contains the \LaTeX\ code to generate the
  label. Backslashes are doubled.
  \index{Label!backslash in}
  \index{Label!double quotes}
\end{itemize}
\index{Label|)}


\subsection{Color Declarations}

\index{Color!declaration}
Old-style color declarations set the fill color, line color, and
text color.
\findex{rgb}
\findex{cmyk}
\begin{verbatim}
  rgb(r, g, b);      cmyk(c, m, y, k);
  rgb(P);            cmyk(P);  // for function-controlled colors
  red(d);  // similarly for other primaries
\end{verbatim}
\index{Color|)}


\section{Creating and Drawing Objects}

Scene elements include geometric objects, coordinate grids, axis
labels, and function plots. A few commands create an object (a point,
line segment, circle, sphere, or plane) that can be used in subsequent
computations, but most drawing commands automatically create, draw,
and discard objects.


\subsection{Geometric Data Structures}

\subsubsection*{Points}
\index{Point|(}
The simplest object in the world, and by far the most common named
data structure, is~\code{P}, an ordered triple of real numbers
(double-precision floats). The function~\code{P(x1,x2,x3)} creates the
point~\teximath{(x_1,x_2,x_3)}. If only two arguments are provided,
\teximath{x_3=0} by default. This convention allows \ePiX\ to treat
2-~and 3-dimensional figures uniformly. The standard basis is
available: \code{E\_1=P(1,0,0)}, etc. 

Depending on context, a \code{P} may represent either a
\emph{location} (point) or a \emph{displacement} (vector). Almost all
\ePiX\ functions treat a~\code{P} as a point. However, algebraic
operators and commands that plot vector fields treat \code{P}
arguments as displacements.

Polar, cylindrical, and spherical coordinate \code{P}~constructors are
sensitive to the current angular mode. 
\index{Angle!mode}
\findex{P}
\findex{polar}
\findex{cis}
\findex{cyl}
\findex{sph}
\begin{verbatim}
  P pt=polar(r,t);   // (r*Cos(t), r*Sin(t), 0)
  P pt=cis(t);       // (Cos(t), Sin(t), 0) = polar(1, t)
  P pt=cyl(r,t,z);   // (r*Cos(t), r*Sin(t), z)
  P pt=sph(r,t,phi); // polar(r,t)*Cos(phi) + (0,0,r*Sin(phi))
\end{verbatim}

Algebraic operations---addition/subtraction, scalar multiplication;
scalar, cross, and componentwise products; orthogonalization---can
be performed on~\code{P}s.  In compound expressions, the binary
operators below should be enclosed in parentheses, and scalars must be
collected at left, \code{P}s at right.
\findex{P.x1}
\findex{P.x2}
\findex{P.x3}
\findex{J}
\begin{verbatim}
  double u=pt.x1();   // first coordinate of pt, etc.
  P(a,b,c)|P(x,y,z);  // scalar product, ax+by+cz
  P(a,b,c)&P(x,y,z);  // componentwise product (ax, by, cz)
  P(a,b,c)*P(x,y,z);  // cross product (bz-cy, cx-az, ay-bx)
  J(p);               // quarter turn about the x3-axis
  p%q;                // orthogonalization, p (mod q)
\end{verbatim}
Explicitly, \code{p\%q}~is the unique vector \code{p+k*q}
perpendicular to~\code{q}.

\code{P}~operations express mathematical relationships, and therefore
imbue a figure with logical structure, making the input file easier to
read, modify, and maintain. Commonly, a file preamble will define a
few named points with hard-coded coordinates, then define additional
points of interest using \code{P}~operators.
\index{Point|)}

\subsubsection*{Complex Numbers}

\index{Complex number|(}
\index{Function!complex}
The \code{Complex} class represents a complex number, specified by its
real and imaginary parts. Arithmetic operators (addition, subtraction,
multiplication, division) are provided, as well as exponential,
logarithmic, circular, and hyperbolic functions (sensitive to angle
units), the latter with \code{C} appended to their names. A complex
number implicitly converts to a point with third coordinate zero. The
sample file \filename{cubic\_cutaway.xp} illustrates usage, including
how to define and plot functions and color-shade Riemann surfaces.
\findex{Complex.re}
\findex{Complex.im}
\findex{Complex.conj}
\findex{Arg}
\findex{norm}
\findex{expC}
\findex{logC}
\findex{sqrtC}
\findex{rootC}
\findex{powC}
\findex{SinC}
\findex{CosC}
\findex{TanC}
\findex{CotC}
\findex{SecC}
\findex{CscC}
\findex{SinhC}
\findex{CoshC}
\findex{TanhC}
\findex{CothC}
\findex{SechC}
\findex{CschC}
\begin{verbatim}
  Complex u, z(4,3);  // u = 0, z = 4 + 3i
  double len(norm(z));// len = 5
  double th(Arg(z));  // th = atan2(3, 4);
  P pt(z);            // pt = (4, 3, 0)
  Complex w(expC(z)); // w = e^{4 + 3i}
  u = LogC(z,k=0);    // u = Log(z) + 2*pi*ki
  u = sqrtC(z,k=0);   // sqrt(z)*expC(pi*ki)
  u = rootC(z,n,k=0); // (z^{1/n})*expC(2*pi*ki/n)
  u = powC(z,n);      // z^n
\end{verbatim}
\index{Complex number|)}


\subsubsection*{Other Geometric Classes}

In addition to~\code{P}, objects of type \code{Circle}, \code{Plane},
\code{Segment}, and \code{Sphere} can be used for Euclidean geometry
constructions. Simple affine operations are supplied for each type, as
is a \code{draw()} function, which represents the object as a path in
the screen.
\findex{shift}
\findex{move\_to}
\findex{scale}
\findex{draw}
\begin{verbatim}
  obj.shift(P arg);    // translate by arg
  obj.move_to(P arg);  // move center to arg
  obj.scale(double c); // scale about center by c
  obj.draw();
\end{verbatim}
A \code{Segment}'s ``center'' is its midpoint. A \code{Plane} has no
center; \code{move\_to} translates the \code{Plane} to pass through
\code{arg}, and \code{scale} has no effect.

\index{Circle}
A \code{Circle} data structure consists of a center, radius, and a
perpendicular unit vector. Three constructors are provided:
\findex{Circle}
\begin{verbatim}
  Circle(center=P(0,0,0), double rad=1, normal=E_3);
  Circle(P center, P point);
  Circle(P p1, P p2, P p3);
\end{verbatim}
The second constructor creates the \code{Circle} parallel to the
\teximath{(x_1,x_2)}~plane, with given center, and radius equal to the distance
between the arguments. (A warning is printed if the second argument
does not lie on the circle, namely, if the arguments do not lie in a
plane parallel to the \teximath{(x_1,x_2)}~plane.) The third returns the
\code{Circle} passing through the given points; the points must not be
collinear.

The data defining a \code{Circle} are recovered with member functions
named \code{center()}, \code{radius()}, and \code{perp()}.

\index{Plane}
\findex{Plane}
A \code{Plane} is specified by a point and normal vector, or by three
non-collinear points. The \code{draw()} function clips the plane and
draws the resulting polygon. Unless the clip box has been set
manually, the clipped polygon's vertices will have large coordinates.

\index{Segment}
\findex{Segment}
A \code{Segment} is constructed from its endpoints. The member
function \code{midpoint()} returns the center.

\index{Sphere}
\findex{Sphere}
A \code{Sphere} is specified by a point and a radius---by default the
origin and unity, or by the center and a point on the sphere. Member
functions \code{center()}~and \code{radius()} return the defining
data.  Capabilities specific to geography and spherical geometry are
described below, pp.~\pageref{non-eucl}ff.

The \code{draw()} function of a \code{Sphere} draws the horizon
visible from the current viewpoint. While this horizon is a circle in
object space, its image in the screen is generally an ellipse.
Antipodal points are not generally mapped to points symmetrically
placed with respect to the center of this ellipse. These effects are
most pronounced when the viewpoint is close to the \code{Sphere} and
the center is not close to the \code{target}.


\subsubsection*{Intersection}

\index{Intersection}

To facilitate geometric computation, \ePiX's \code{Circle},
\code{Plane}, \code{Segment}, and \code{Sphere} classes can be
intersected with the \code{*}~operator.
\texicomment
Table~\ref{table-intersect}
lists the return types for each pair of arguments.  Intersection is
commutative, so only the top half of the table is shown. For purposes
of intersection, a \code{Segment} is extended into a line.
The sample file \filename{pascal.xp} gives typical applications of
objects and intersection.
\findex{intersection}
\begin{table}[hbt]
\begin{center}
  \begin{tabular}{c|cccc}
    \code{*} & 
    \code{Segment} & \code{Circle} & \code{Plane} & \code{Sphere} \\
    \hline
    \code{Segment} & \code{P} & \code{Segment} & \code{P} & \code{Segment} \\
    \code{Circle}  &    & \code{Segment} & \code{Segment} & \code{Segment} \\
    \code{Plane}   &    &    &   \code{Segment} & \code{Circle} \\
    \code{Sphere}  &    &    &    & \code{Circle}
  \end{tabular}
  \caption{Object intersection types.}
  \label{table-intersect}
\end{center}
\end{table}
\texicommentend

A \code{Circle} has a center, radius, and unit normal; a \code{Plane}
has a distinguished point and unit normal; a \code{Segment} has two
endpoints; a \code{Sphere} has a center and radius. An object is
\emph{malformed} if these conditions are not met.  The constructors
return well-formed objects with two exceptions: \code{Circle} and
\code{Plane} create malformed objects if called with three collinear
points.  The operator~\code{*} returns a malformed object if either
argument is malformed, or if the operands are disjoint, tangent, or
coincident.  Malformedness is benign: Calling \code{draw()} on a
malformed object does nothing.


\subsubsection*{Orthonormal Frames}

A \code{frame} comprises three mutually perpendicular unit
vectors. The constructor takes three vectors. The \code{frame}'s third
vector~\teximath{e_3} is positively proportional to~\code{v3}, the
second vector~\teximath{e_2} is positively proportional to
\code{v2\%v3}, and the first is the cross product,
\teximath{e_1=e_2\times e_3}. Thus, a \code{frame} is right-handed,
and does not depend on~\code{v1}. 

The elements of a \code{frame} are named \code{sea}, \code{sky}, and
\code{eye}, just as for the \code{camera}. A \code{frame} can be
rotated through an arbitrary angle about any of its elements.
\findex{frame}
\begin{verbatim}
  frame();        // the standard basis {E_1, E_2, E_3}
  frame fr(v1, v2, v3); // orthonormalize {v1, v2, v3}
  fr.sea();       // the first element of fr, etc.
  fr.rot1(theta); // rotate fr through theta about sea, etc.
\end{verbatim}


\subsection{Path-Like Elements}

Basic path-like objects are drawn with named commands. Arguments of
polygon commands are endpoints/vertices. Except for \code{line}~and
\code{Line}, the following are subject to filling.
\findex{line}
\findex{Line}
\findex{triangle}
\findex{rect}
\findex{quad}
\findex{circle}
\begin{verbatim}
  line(P p1, P p2, [double expand]);
  Line(p1, p2);  // draw line through p1, p2 (crop required)
  triangle(P p1, P p2, P p3);
  rect(P p1, P p2);
  quad(P p1, P p2, P p3, P p4); // quadrilateral
  circle(ctr=P(0,0,0), rad=1, normal=E_3);
  circle(ctr, pt);
  circle(pt1, pt2, pt3);
\end{verbatim}
The optional \code{line} argument is an expansion parameter:
\code{line(p1,p2,t);} draws a segment centered at the midpoint of
\code{p1}~and \code{p2}, with length scaled
by~\teximath{2^{t/100}}. (Setting \teximath{t=100} doubles the length,
while \teximath{t=-100} halves the length.) The arguments of
\code{rect()} must lie in a plane parallel to a coordinate plane. The
arguments to \code{circle} commands are the same as for \code{Circle}
constructors.

Quadratic and cubic splines are described by their control points.  A
list of~\code{P} is drawn as a ``natural'' spline (the \teximath{C^2}
piecewise cubic curve with vanishing second derivatives at the
endpoints); the number of points per cubic segment must be specified.
Circular and elliptical arcs are given by center, a basis, angular
range, and an optional number of intervals.
\findex{spline}
\findex{arc}
\findex{ellipse}
\begin{verbatim}
  spline(P p1, P p2, P p3, [int n]);       // quadratic
  spline(P p1, P p2, P p3, P p4, [int n]); // cubic
  spline(vector<P>, int n);                // natural spline

  arc(P ctr, rad, t_min, t_max);  // parallel to (x1,x2)-plane
  ellipse(P ctr, P v1, P v2);     // in plane spanned by v1, v2
  ellipse(P ctr, P v1, P v2, t_min, t_max, [int n]);
\end{verbatim}
\texicomment
Mathematically, these commands draw parametric paths
\begin{align*}
  \text{Spline:}\quad&(1-t)^2 p_1 + 2(1-t)t p_2 + t^2 p_3,
    && t\in[0,1] \\
  \text{Spline:}\quad&(1-t)^3 p_1 + 3(1-t)^2t p_2 + 3(1-t)t^2 p_3 + t^3p_4,
    && t\in[0,1] \\
  \text{Arc:}\quad&\mathrm{ctr} + (\cos t)(\mathrm{rad},0,0)
                + (\sin t)(0,\mathrm{rad},0),
    && t\in[t_\mathrm{min}, t_\mathrm{max}] \\
  \text{Ellipse:}\quad&\mathrm{ctr} + (\cos t)v_1 + (\sin t)v_2, 
    && t\in[t_\mathrm{min}, t_\mathrm{max}].
\end{align*}
\texicommentend
If parameter bounds are omitted in an \code{ellipse} command, the
entire ellipse is drawn.  When the angular range subtends one or more
full turns in an \code{arc}~or \code{ellipse} the curve is subject to
filling.

Commands for planar (half-)ellipses remain from \ePiX's early days:
\findex{ellipse}
\findex{ellipse\_left}
\findex{ellipse\_right}
\findex{ellipse\_top}
\findex{ellipse\_bottom}
\begin{verbatim}
  ellipse(P ctr, P radius);
  ellipse_left(P ctr, P radius);
  ellipse_right(P ctr, P radius);
  ellipse_top(P ctr, P radius);
  ellipse_bottom(P ctr, P radius);
\end{verbatim}
If \code{radius} is \code{P(a,b)}, these commands draw all or half of
the ellipse with given center in the \teximath{(x_1,x_2)} plane, axes
parallel to the coordinate axes, and axis lengths \teximath{2a}~and
\teximath{2b}.

Two commands are available to mark off right angles or a subtended
angle. Each accepts a spatial location and two non-zero directions,
and draws a scene element in the plane spanned by the vectors.
\findex{right\_angle}
\findex{arc\_measure}
\begin{verbatim}
  right_angle(P loc, P v1, P v2, scale=8);
  arc_measure(P loc, P v1, P v2, scale=8);
  arc_measure(P loc, P v1, P v2, offset, text, align, scale=8);
\end{verbatim}
The \code{right\_angle} command does not check its arguments for
perpendicularity. The \code{arc\_measure} commands mark the
\emph{small} angle subtended by the directions; the label form places
a label at the midpoint of the arc. The (optional) \code{scale}
argument is the true size in~\code{pt} of the marker.


\begin{comment}
\subsubsection*{Recursive Fractal Curves}  

\index{Path!fractal}
Consider a path made up of equal-length segments that can point at any
angle of the form~\teximath{2\pi k/n} radians, for \teximath{0\leq
k<n}, like spokes on a wheel. A path is specified by a finite sequence
of integers, taken modulo~\teximath{n}. For example, if
\teximath{n=6}, then the sequence \teximath{0, 1, -1, 0} corresponds
to the ASCII path \verb+_/\_+.  \ePiX's fractal approximation starts
with such a ``seed'' then recursively (up to a specified depth)
replaces each segment with a scaled and rotated copy of the seed,
Figure~\ref{fig-fractal}. The seed above generates the standard
von~Koch snowflake fractal. In code: 
\findex{fractal}
\begin{verbatim}
  const int seed[] = {6, 4, 0, 1, -1, 0};
  fractal(P(a,b), P(c,d), depth, seed);
\end{verbatim}
The first entry of \code{seed[]} (here~6) is the number of
``spokes''~\teximath{n}, the second~(4) is the number of terms in the
seed, and the remaining entries are the seed proper. The final path
joins \teximath{(a,b)}~to \teximath{(c,d)}. The number of segments in
the final path grows exponentially in the depth, so depths larger than
5~or 6 are likely to exceed the capabilities of \LaTeX\ and/or
PostScript.
\begin{figure}[hbt]
  \begin{center}
  \input{koch.eepic}
  \end{center}
  \caption{Successive iterations of \{\code{4,8,0,1,0,3,3,0,1,0}\}}
  \label{fig-fractal}
\end{figure}
\end{comment}


\subsubsection*{Arrows}

\index{Arrow}
Line segments, splines, and arcs can be drawn with arrowheads at one
end.  In profile, an arrowhead's width is \code{3pt}, and its height
is 5.5~times the width.  The actual printed height depends on the
head's orientation with respect to the camera.  By default, an
arrowhead is a filled triangle. The shape and size are adjusted with
declarations:
\findex{arrow\_width}
\findex{arrow\_ratio}
\findex{arrow\_inset}
\begin{verbatim}
  arrow_width(w=3);   // width in pt
  arrow_ratio(r=5.5); // height-to-width
  arrow_inset(c=0);   // base indent as frac of ht
\end{verbatim}
\input{arrow.eepic} Inset\teximath{=0}\qquad
\input{arrow25.eepic} Inset\teximath{=0.25}\qquad
\input{arrow50.eepic} Inset\teximath{=0.5} \\
The \code{inset} must be between \teximath{-1}~and \teximath{1}. Each
\code{arrow} command accepts an optional \code{scale} argument, which
scales the arrowhead.
\findex{arrow}
\begin{verbatim}
  arrow(P tail, P tip, [scale]);
  arrow(P p1, P p2, P p3, [scale]);   // spline arrows
  arrow(P p1, P p2, P p3, P p4, [scale]);
  arrow(P ctr, P v1, P v2, t_min, t_max, [scale]); // ellipse
\end{verbatim}
A few ``special-purpose'' commands are supplied:
\findex{dart}
\findex{aarrow}
\findex{arc\_arrow}
\begin{verbatim}
  dart (P p1, P p2);  // same as arrow(p1, p2, 0.5);
  aarrow(P p1, P p2); // double-headed arrow <--->
  arc_arrow(ctr, rad, t_min, t_max, [scale]);
\end{verbatim}
If an \code{arc\_arrow} is too short, nothing is drawn.


\subsection{Coordinate Axes and Labels} 

\index{Label!axis|(}
\index{Axes!labeling}
\index{Axes|(}
\ePiX\ provides an \code{axis} class for coordinate axes. Labels are
generated automatically in a variety of styles: decimal, scientific
notation, fraction, and trigonometric fraction. Logarithmic axes and
labels are available. Axis and labeling commands from Version~1.0 have
been retained.


\subsubsection*{The \code{axis} Class}

An \code{axis} consists of a line segment divided into a specified
number of equal-length subintervals by ``major'' (long) tick marks.
Each subinterval may be divided further by ``minor'' (short) tick
marks.  Minor ticks may be equally-spaced (``Cartesian'') or
logarithmically placed. Finally, a label is written at each major tick
mark with specified offset and alignment. Labels are generated
automatically from the endpoints, so the line should be parallel to a
coordinate axis. Label attributes (masking, borders, font size and
face, rotation angle) are determined by the current label style, not
by the \code{axis}.

The command
\findex{axis}
\begin{verbatim}
  axis(P tail, P head, int n, P offset=P(0,0), [align = none]);
\end{verbatim}
creates an axis joining \code{tail} to \code{head}, divided into
\code{n}~segments of equal length, with a major tick mark and label at
each division point and endpoint. The \code{offset} and \code{align}
arguments have the same meaning as for ordinary labels.

The number and length of minor ticks, and the alignment of ticks, are
controlled with member functions:
\findex{axis.subdivide}
\findex{axis.tick\_ratio}
\findex{axis.align}
\begin{verbatim}
  axis Ax(P(a,c), P(b,c), n);
  Ax.subdivide(n);  // put n-1 minor ticks in each axis segment
  Ax.tick_ratio(r); // minor length = r*major length
  Ax.align(AL=c);   // align all ticks; AL = t, b, l, r, or c
\end{verbatim}
By default, major ticks are \code{6pt}~long and twice the length of
minor ticks. The global declaration \code{tick\_size(len)} sets the
major length, subject to ``reasonable stylistic limits''. For visual
consistency, tick lengths should not be changed casually. 

\index{Axes!labeling}
Labels on an \code{axis} may be drawn in several styles, selected with
member functions:
\findex{axis.dec}
\findex{axis.frac}
\findex{axis.trig}
\findex{axis.sci}
\findex{axis.unmark}
\findex{axis.precision}
\findex{axis.align\_labels}
\begin{verbatim}
  Ax.dec();  // decimals (default)
  Ax.frac(); // fractions, e.g. 0.5 -> \frac{1}{2}
  Ax.trig(); // fractional multiples of \pi
  Ax.sci();  // scientific notation, k\times base^N

  Ax.unmark(double); // remove label at selected location
  Ax.precision(p);   // set number of digits for decimal labels
  Ax.align_labels(AL); // re-align labels
\end{verbatim}
Arbitrary textual labels depending on one coordinate are obtained by
writing a \code{string}-valued function of~\code{double} and
``registering'' it:
\findex{axis.label\_rep}
\begin{verbatim}
  // f represents x as a string in given precision and base
  std::string f(double x, unsigned int prec, unsigned int base);
  Ax.label_rep(f);
\end{verbatim}
For example, textual tags can be printed instead of numeric labels.

By default, an \code{axis} is Cartesian. The member function
\code{log(int b=10)} converts an \code{axis} to ``log mode'' with
specified base; this affects both tick marks and labeling. If \code{b}
is at least~3, minor tick marks appropriate for a logarithmic axis
base~\code{b} are drawn. Second, labels are written in decimal or
scientific notation appropriately for a logarithmic axis; that is,
``\teximath{k\times b^N}'' (or its decimal value) is written at
location \teximath{N+\log_b k}.

Minor ticks of a log \code{axis} may be labeled individually; again,
this is controlled with member functions:
\findex{axis.tag}
\findex{axis.tag235}
\findex{axis.tags}
\findex{axis.untag}
\begin{verbatim}
  Ax.tag(d);   // labels at N+log_b(d)
  Ax.tag235(); // tag 2, 3, and 5 if b=10
  Ax.tags();   // tag 1, ..., b-1
  Ax.untag();  // remove all tags, including 1
  Ax.untag(double);  // remove one tag, e.g. Ax.untag(9);
\end{verbatim}

For convenience, an \code{axis} object along an edge of the active
\code{screen} can be created with a named command:
\findex{top\_axis}
\findex{bottom\_axis}
\findex{left\_axis}
\findex{right\_axis}
\begin{verbatim}
  top_axis(n, offset, align);
  bottom_axis(n, offset, align);
  left_axis(n, offset, align);
  right_axis(n, offset, align);
\end{verbatim}
The ticks automatically point into the bounding box.

The functions above create objects or set \code{axis} attributes, but
write no output. Tick marks and labels can be printed separately, or
at once:
\findex{axis.draw}
\findex{axis.draw\_ticks}
\findex{axis.draw\_labels}
\begin{verbatim}
  Ax.draw();        // axis, tick marks, and labels
  Ax.draw_ticks();  // axis and ticks only
  Ax.draw_labels();
\end{verbatim}


\subsubsection*{Other Axis-Drawing Commands}

\index{Axes!logarithmic}
\ePiX\ formerly supplied commands for drawing simple axes and their
labels. These commands have been kept for compatibility.  Horizontal
axes are generated with
\findex{h\_axis}
\findex{h\_log\_axis}
\findex{v\_axis}
\findex{v\_log\_axis}
\begin{verbatim}
  h_axis(p1, p2, n, align=c);    // n subintervals (n+1 ticks)
  h_log_axis(p1, p2, n, align=c, base=10);
\end{verbatim}
For vertical axes, use \code{v\_axis}.  The style of tick mark is
appropriate for an axis of the given type.  Horizontal axis tick marks
may be aligned \code{t}~(above the axis) or~\code{b}
(below). Similarly, vertical axis ticks may be aligned \code{r}~or
\code{l}.

The endpoint arguments of a coordinate axis may be omitted, in which
case they default to \teximath{p_1=(x_\mathrm{min},0)} and
\teximath{p_2=(x_\mathrm{max},0)} for a horizontal axis, or to
\teximath{p_1=(0,y_\mathrm{min})} and
\teximath{p_2=(0,y_\mathrm{max})} for a vertical axis. If the bounding
box has integer width and/or height, omitting the number of points
draws tick marks one unit apart.

Labels for a horizontal Cartesian or logarithmic axis are generated
with the commands
\findex{h\_axis\_labels}
\findex{h\_axis\_masklabels}
\findex{h\_axis\_log\_labels}
\findex{h\_axis\_log\_masklabels}
\findex{v\_axis\_labels}
\findex{v\_axis\_masklabels}
\findex{v\_axis\_log\_labels}
\findex{v\_axis\_log\_masklabels}
\begin{verbatim}
  h_axis_labels(P p1, P p2, int n, P offset, [align]);
  h_axis_masklabels(p1, p2, n, offset, [align]);
  h_axis_log_labels(p1, p2, [n], offset, [align], base=10);
  h_axis_log_masklabels(p1, p2, [n], offset, [align], base=10);
\end{verbatim}
Labels for a vertical axis are generated with \code{v\_axis\_labels},
etc. The labels are automatically generated to match their horizontal
location.  The first puts \teximath{(n+1)}~evenly-spaced labels on the segment
joining \code{p1}~and \code{p2}. As with ordinary labels, the
\code{offset} is in~\code{pt}, and the optional \LaTeX-style alignment
option positions the labels using their corners.  The second command
draws masked labels according to the current label masking
attributes. The third writes labels in exponential notation, using the
Cartesian coordinate as exponent.

As for coordinate axes, the initial and final points may be omitted in
an \code{axis\_[mask]labels} command, with the same defaults.  The
\code{offset} and number of labels must always be specified.
\index{Label!axis|)}


\subsubsection*{Broken Axes}
\index{Axes!broken}

Broken axes are best drawn using page layout, especially if axis
labels are to be drawn. To accomplish the task, create a screen for
each piece of axis, using appropriate Cartesian coordinates for the
corners, then \code{inset} the screens so the axis pieces are nearly
end to end. A zig-zag glyph signifies the break:
\findex{axis\_break}
\begin{verbatim}
  axis_break(P, P, scale=12);
\end{verbatim}
The \code{P} arguments are the screen coordinates of the ends to be
joined; the optional third argument is the true height and width
in~\code{pt} of the glyph. The sample file \filename{coord\_tricks.xp}
uses layout and axis breaks.


\subsubsection*{Coordinate Grids}

Cartesian grids fill a coordinate rectangle, and have a specified
number of lines in each direction. A polar grid has specified radius,
and numbers of rings and sectors.
\findex{grid}
\findex{polar\_grid}
\begin{verbatim}
  grid(n1, n2);         // fills the bounding box
  grid(p1, p2, n1, n2); // fills the box with corners p1, p2
  polar_grid(r, n1, n2);
\end{verbatim}

Log and semi-log grids are drawn with analogous syntax:
\findex{log\_grid}
\findex{log1\_grid}
\findex{log2\_grid}
\begin{verbatim}
  log_grid (p1, p2, n1, n2, [base1], [base2]); // log-log
  log1_grid(p1, p2, n1, n2, [base]);           // log-lin
  log2_grid(p1, p2, n1, n2, [base]);           // lin-log
\end{verbatim}
The~\teximath{n_i} arguments dictate the number of squares, namely the number
of orders of magnitude spanned in the logarithmic direction(s).  The
(optional) base arguments default to~10, and control the number of
lines drawn per order of magnitude. As with Cartesian grids, the
corners are optional; if omitted, the grid fills the bounding box.

\index{Graph paper}
Graph paper may be created by superimposing grids:
\texicomment
\begin{center}
  \begin{minipage}[b]{2.5in}
    \begin{footnotesize}
\texicommentend
\findex{pen}
\findex{grid}
\begin{verbatim}
  pen(0.25);
  grid(10*xsize(), 10*ysize());
  pen(0.5);
  grid(2*xsize(), 2*ysize());
  pen(1);
  grid(xsize(), ysize());
\end{verbatim}
\texicomment
    \end{footnotesize}
  \end{minipage}
  \qquad
  \input{graphpaper.eepic}
\end{center}
\index{Axes|)}
\texicommentend

\subsection{The Path Class}

\index{Path!class}
A \code{path} data structure is an ordered list of points that can be
cropped, clipped, mapped, concatenated, and drawn.  The first four
constructors mirror polygon-drawing commands above.  Function graphs
and parametric paths are built from a real- or \code{P}-valued
function~\code{f} of one variable.
\findex{path}
\begin{verbatim}
  path(p1, p2, expand=0);     // line (endpoints)
  path(p1, p2, p3, [n]);      // quadratic spline
  path(p1, p2, p3, p4, [n]);  // cubic spline
  path(p1, v1, v2, t_min, t_max, [n]); // ellipse
  path(f, t_min, t_max, [n]);
\end{verbatim}

The member function \code{pt()} accepts a~\code{P} or three (or two)
doubles and appends the specified point to a \code{path}.  This
snippet creates a regular \teximath{n}-gon:
\begin{verbatim}
  path ngon; // declare new path
  for (int i=0; i<=n; ++i)
    ngon.pt(cis(i*full_turn()/n)); // works in all angle modes
\end{verbatim}
Compound paths may be built by concatenation.  If \code{path1}~and
\code{path2} share an endpoint, the commands
\begin{verbatim}
  path1 += path2;
  path1 -= path2;
\end{verbatim}
replace \code{path1} with the result of traversing \code{path1}
``forward'', then following \code{path2} in the forward or reverse
direction (respectively). For expected results, the first
(or last) point on~\code{path2} should be the last point
of~\code{path1}. The notation suggests 1-dimensional homology
chains. The sample file \filename{contour.xp} illustrates path
creation and manipulation.

A \code{path} is a data structure, and must be drawn explicitly to
create visible output. By default a \code{path} is not a closed loop
(even if the first and last points are the same), and is not filled
when drawn. Member functions perform these tasks. Continuing the
\teximath{n}-gon snippet above,
\findex{path.close}
\findex{path.fill}
\findex{path.draw}
\begin{verbatim}
  ngon.close(); // mark path as closed
  ngon.fill();  // draw filled region if filling is active
  ngon.draw();  // print to the screen
\end{verbatim}
\code{path::close()} adds a closing edge if necessary. Once a
\code{path} is closed, no more points can be added. A closed path
clips and crops differently than an open \code{path} with the same
data. \code{path::fill()} has no effect on an unclosed \code{path}.


\subsection{Function Plotting}

\index{Plotting|(}
The noun ``map'' refers to a \CXX\ function that accepts one or
more \code{double} arguments and returns a \code{double}~or
a~\code{P}. Mathematically, a map can be depicted in two ways: as a
graph (which retains information about the domain), or as a
parametrized curve or surface (which discards domain
information). \ePiX\ assumes that \code{double}-valued maps are
graphed and \code{P}-valued maps are drawn parametrically. Either sort
of depiction is called a ``plot''.  \ePiX\ plots are either ``wire
mesh'', produced by a \code{plot} command, or ``shaded'', produced by
a \code{surface} command.


\subsubsection*{Basic Plotting}

For the moment, ``function'' means ``function of one variable''
(precisely, a \code{double}-valued function of a \code{double}
variable). A function graph depends on the domain and the number of
points to use. Each of the commands
\findex{plot}
\findex{polarplot}
\findex{shadeplot}
\begin{verbatim}
  plot(f, t_min, t_max, n);
  polarplot(f, t_min, t_max, n);
  shadeplot(f, t_min, t_max, n);
\end{verbatim}
graphs the function~\code{f} on the interval \code{[t\_min, t\_max]}
by dividing the interval into \code{n}~subintervals of equal length.
The first gives a Cartesian plot, the second a polar plot with bounds
in current angular units, the third shades the region between the
graph and the horizontal axis. If two functions are given to
\code{shadeplot}, the region between their graphs is shaded.


\subsubsection*{Domains and Wiremesh Plots}

\index{Domain|(}
An \ePiX\ \code{domain} is a coordinate box of dimension one, two, or
three, specified by a pair of opposite corners and two \emph{meshes}
(``coarse'' and ``fine'', respectively) which specify the amount of
data to be plotted. Plotting is explained in detail below.
\findex{domain}
\begin{verbatim}
  // [a1,a2] x [b1,b2]: n1 x n2 rectangles, m1 x m2 intervals
  domain R2(P(a1,b1), P(a2,b2), mesh(n1,n2), [mesh(m1,m2)]);

  // [a1,a2] x [b1,b2] x [c1,c2] divided analogously
  domain R3(P(a1,b1,c1), P(a2,b2,c2),
            mesh(n1,n2,n3), [mesh(m1,m2,m3)]);
\end{verbatim}
If unspecified, the fine mesh is the same as the coarse mesh. For
expected behavior, the coarse mesh should ``divide'' the fine mesh, in
that \teximath{m_i}~should be a (usually small) integer multiple
of~\teximath{n_i} for each~\teximath{i}.

A \code{domain} may be \emph{resized} in any coordinate for which the
thickness is positive, and can be \emph{sliced} by setting one
variable to a constant. The result of slicing is a \code{domain} whose
dimension is one smaller than the original.  Finally,
``\code{slices}'' operators return the list of \code{domain}s obtained
by setting one variable to evenly-spaced constants.  By default, the
number of slices is specified by the coarse mesh. An optional argument
specifies the number of slices. This argument need not be related to
the coarse mesh.
\findex{domain.resize}
\findex{domain.slice}
\findex{domain.slices}
\begin{verbatim}
  R2.resize2(a,b); // [a1,b1] x [a,b]
  R2.slice1(t);    // set x1 = t
  R3.slices3([n]); // (n+1) domains with x3 = const
\end{verbatim}
\begin{center}
  \input{slices.eepic}
\end{center}
\smallskip
When possible, resizing preserves grid square sizes. Generally,
though, integer truncation occurs: If \code{R=domain(P(0,0), P(1,1),
mesh(10,6))}, then \code{R.resize1(0,0.25)} is the rectangle
\teximath{[0,0.25]\times[0,1]} subdivided into \teximath{2\times6}
subrectangles, since \teximath{10\div4=2} in integer arithmetic. For
expected behavior, choose mesh sizes to avoid integer truncation.


The arguments of a \code{plot} command are a map, followed by either a
domain or its logical equivalent.
\findex{plot}
\begin{verbatim}
  double f(double t) { return t*t; }
  P F(double u, double v) { return P(u, v, exp(u)*Sin(v)); }
  P G(double u, double v, double w) { return P(v*w, u*w, u*v); }

  plot(f, a, b, n);       // f:[a,b] -> R, using n intervals
  plot(F, R2);            // graph of exp(u)*Sin(v)
  plot(G, R3.slice2(0.5); // G: R^3 -> R^3 restricted to y=0.5
\end{verbatim}
By (compiler-enforced) convention, \code{plot} commands involving a
\code{P}-valued map accept a \code{domain} argument, as in the second
and third commands above. To plot a \code{double}-valued function, by
contrast, supply the logical equivalent of a \code{domain}, usually
the endpoints and the number of intervals, as in the first \code{plot}
command above.

Resizing and slicing allow a map~\code{F} to be plotted selectively
over parts of its domain. This can be used to emphasize parts of the
image, layer scene elements, patch surfaces together, and so forth.
Resize and slice(s) commands may be used directly in a \code{plot}
command:
\findex{plot}
\begin{verbatim}
  plot(F, R2.resize1(0,0.5));
  plot(F, R2.slices1());
\end{verbatim}


\subsubsection*{Meshes and Plotting}

The \code{P}~arguments of a \code{domain} are a pair of opposite
corners. The first \code{mesh} argument, the \emph{coarse} mesh,
specifies the number of subdivisions in each coordinate direction.
The second \code{mesh}, the \emph{fine} mesh, determines the number of
points used in each direction when plotting.

\begin{comment}
Suppose the \code{domain}~\code{R2} is, as above, the rectangle
\teximath{[0,3]\times[-1,2]}. The argument \code{mesh(12,6)} divides the
rectangle a \teximath{12\times6}~array of subrectangles known as the
\emph{coarse mesh}. A \code{P}-valued map~\code{F} of two variables is
plotted over the edges of the coarse mesh. Rather than drawing
quadrilaterals, however, \ePiX\ draws the image curves at the
resolution specified by the argument \code{mesh(60,60)}, the
\emph{fine mesh}. In this example, the curves in each direction are
drawn using 60~line segments.
\end{comment}

Separating the roles of coarse and fine meshes allows a plot to
conform closely to a surface without using a fine grid of curves.
Both parts of Figure~\ref{fig-plot} are drawn with a
\teximath{6\times20} coarse mesh. In the first picture, the fine mesh
is also \teximath{6\times20}, while in the second, the fine mesh is
\teximath{12\times60}.

\begin{figure}[hbt]
  \begin{center}
    \input{plot_a.eepic}\hspace*{0.5in}
    \input{plot_b.eepic}
    \caption{Coarse and fine meshes.}
    \label{fig-plot}
  \end{center}
\end{figure}

The coarse mesh is significant only for domains of dimension at
least~2. The coarse mesh's size determines the number of curves or
surfaces plotted \emph{perpendicularly} to a coordinate direction,
while the fine mesh's size determines the number of segments used
\emph{along} that direction.  For predictable results, the fine mesh
should be a small multiple of the coarse mesh.

Plotting works analogously for 3-dimensional \code{domain}s and maps
depending on three variables: The ``one-dimensional skeleton'' of the
\code{domain}'s image is drawn.  A \code{P}-valued map of
three variables can be plotted over a 1-~or 2-dimensional
\code{domain}. (The effect may be unexpected unless the domain arises
by slicing, however.) A map depending on one or two variables cannot
be plotted over a 3-dimensional \code{domain}.
\index{Domain|)}


\subsubsection*{Shaded Surfaces}

\index{Plotting!surface|(}

\ePiX's shaded surface plotting implements a degree of hidden surface
removal. The algorithm breaks a surface into mesh fragments, sorts
them in decreasing (approximate) distance to the camera, and prints
them.  If filling is active, a mesh fragment is shaded according to
the angle between the normal vector and the direction to the camera,
simulating constant ambient lighting; otherwise, the current fill
color is used.  This technique works fairly well for surfaces without
intersection, and even acceptably handles intersecting surfaces for
which mesh elements intersect only along boundaries.

The syntax of a \code{surface} command is identical to a \code{plot}
command when only a single surface is drawn. For example,
\findex{surface}
\begin{verbatim}
  surface(F, R, cull=0);
\end{verbatim}
plots the \code{P}-valued function~\code{F} over the domain~\code{R},
the shaded equivalent of the corresponding \code{plot} command. The
optional \code{cull} argument removes elements that point toward
(\code{cull=-1}) or away from (\code{cull=1}) the camera. Naturally,
orientation of mesh elements depends on the parametrization~\code{F},
not merely on the surface. Culling reduces the output file size, but
is useful mostly for closed, convex surfaces.

There are special commands for surfaces of rotation; each accepts a
final \code{cull} argument.
\findex{surface\_rev}
\begin{verbatim}
  surface_rev(f, t_min, t_max, n_lats, n_longs);
  surface_rev(f, g, t_min, t_max, n_lats, n_longs=24);
  surface_rev(f, g, R, frame coords);
\end{verbatim}
The first revolves the graph of~\teximath{f} about the
\teximath{x}-axis, the second uses the parametric curve
\teximath{t\mapsto\bigl(f(t),g(t)\bigr)} as profile. In each case, the
parameter interval \verb+[t_min,t_max]+ is divided into \verb+n_lats+
equal-length subintervals, \verb+n_longs+ copies of the profile curve
are drawn, and the complete surface (one full turn) is drawn.

The third form uses a \code{domain} to control the range of
longitudes, and draws a surface of rotation in the Cartesian
coordinate system defined by the orthonormal basis \code{coords}, by
default the standard basis. The arguments \code{f}~and \code{g} define
a parametric curve in the plane spanned by the first two elements
of~\code{coords}, and the first element is the axis of rotation.

As in wire mesh plotting, the fine mesh is used to draw the boundaries
of surface patches; this tends to make surfaces look smoother for
modest-sized coarse meshes. If the coarse mesh is too coarse, however,
two visually undesirable effects can occur. First, adjacent regions of
the surface may be shaded very differently, since shading is constant
over patches defined by the coarse mesh. Second, a patch nearly
tangent to a line of sight may be drawn badly if the patch bends back
on itself, since the boundary of the \emph{patch} is drawn, not the
visible edge of the mathematical surface. See
\filename{samples/artifacts.xp}.


\subsubsection*{Multiple Domains and/or Maps}

A scene containing two or more shaded surfaces cannot generally be
built up one surface at a time. Instead, multiple surfaces must be
assembled into a single data structure before they can be drawn.
Multiple surfaces are built from one or more maps and one or more
2-dimensional \code{domain}s. In the code snippets below, \code{F}~and
\code{G} are \code{P}-valued functions of 3~variables, and~\code{R} is
a 3-dimensional \code{domain}.

To plot the images of several \code{domain}s under a \emph{single
map}, assemble the \code{domain}s into a list if necessary, then issue
a \code{surface} command:
\findex{surface}
\findex{domain\_list}
\findex{domain\_list.add}
\begin{verbatim}
  surface(F, R.slices3(), cull=0);

  domain_list DL(R.slice1(0));  // build domain list
  DL.add(R.slice2(0.5));        // add a domain, etc.
  surface(G, DL, cull=0);       // draw
\end{verbatim}
For multiple maps, \ePiX\ provides the \code{scenery}
class. Conceptually, \code{scenery} is an agglomeration of shaded
surfaces, built one surface at a time from maps and 2-dimensional
\code{domain}s. The \code{add} function accepts two arguments---a map,
and either a \code{domain} or a list of \code{domain}s---and
contributes its data to the \code{scenery} rather than plotting
immediately. Completed \code{scenery} is drawn manually.
\findex{scenery}
\findex{scenery.add}
\findex{scenery.draw}
\begin{verbatim}
  scenery S(F, R.slice3(0.25)); // S contains one surface
  S.add(F, R.slice2(0));        // S contains two surfaces
  S.add(G, R.slices1(3));       // S contains six surfaces
  S.draw(cull=0);
\end{verbatim}
Complete examples are included in the \filename{samples} directory:
\filename{spherical.xp}~and \filename{minkowski.xp}.

In principle, a scene may contain arbitrarily many surfaces. However,
figures that contain many objects tend to tax \LaTeX's internal
stacks. Frequent color changes exacerbate the problem.  Even if you
use \filename{hugelatex} (or increase \LaTeX's memory), a figure
containing more than a few thousand mesh elements is unlikely to
compile. At moderate resolution, a surface can easily contain
1000~patches. Each shell script has a command-line option to invoke
\filename{hugelatex}; your mileage may vary.


\subsubsection*{User-Specified Color Shading}

\index{Color!surface plotting}
By default, a \code{surface} or \code{scenery} is colored according to
the current fill color. For finer control, each \code{surface},
\code{surface\_rev}, and \code{scenery} command accepts an optional
position-dependent \code{color} specification.
\findex{surface}
\findex{surface\_rev}
\findex{scenery}
\findex{scenery.add}
\begin{verbatim}
  surface(F, R, color, cull=0);
  surface_rev(f, [g], t_min, t_max, n_lats, n_longs, color);
  surface_rev(f, g, R, color, [coords]);

  scenery S(F, R, color);
  S.add(F, R, color);
\end{verbatim}
The \code{color} argument is a \code{P}-valued function of two or
three \code{double}s whose output is interpreted as a set of
\code{RGB}~densities. If \code{color} takes two arguments, they are
\code{domain} coordinates, and the surface is colored according to
parameter values. If \code{color} takes three arguments, they are
Cartesian coordinates, and the surface is colored according to spatial
location.  Please see the sample files
\filename{surface\_shade.xp}~and \filename{S2\_harmonics.xp}. 
\index{Plotting!surface|)}


\subsection{Calculus Plotting} 
\index{Plotting!calculus|(}

\ePiX\ provides high-level commands for plotting derivatives and
definite integrals, Riemann sums, tangent lines, slope- and vector
fields, and solutions of planar and spatial systems of differential
equations.

In this section, \code{f}~and \code{g} are \code{double}-valued
functions of one variable.

\subsubsection*{Utility Functions}

\findex{sup}
\findex{inf}
\findex{newton}
\begin{verbatim}
  sup(f, a, b);     // max/min of f on [a,b]
  inf(f, a, b);
  newton(f, g, x0); // find approximate crossing point
\end{verbatim}
Newton's method returns the crossing point of the given functions,
starting from the specified seed, which should be reasonably close to
the expected solution. If a critical point is hit or 5~iterations
pass, a warning is issued and the current result (probably incorrect)
is returned. The second function~\teximath{g} defaults to the zero function if
omitted.

\subsubsection*{Derivatives and Integrals}

The classes \code{Deriv}~and \code{Integral} are used to calculate
values of derivatives and integrals, and to plot these functions.
\findex{Deriv}
\findex{Deriv.eval}
\findex{Deriv.left}
\findex{Deriv.right}
\findex{Integral}
\findex{Integral.eval}
\begin{verbatim}
  Deriv df(f); // function object: df(x) = f'(x)
  df.eval(t);  // return f'(t)
  df.left(t);  // deriv from left at t:  (f(t)-f(t-dt))/dt
  df.right(t); // deriv from right at t: (f(t+dt)-f(t))/dt

  Integral prim(f,a); // function object: prim(x) = int_a^x f
  prim.eval(b);  // numerical integral of f over [a,b]
  double val(Integral(f).eval(1));   // val = \int_0^1 f
\end{verbatim}
The lower limit on an integral is~0 by default. \code{Deriv}s and
\code{Integral}s can be used directly in a \code{plot} command:
\findex{plot}
\findex{riemann\_sum}
\begin{verbatim}
  plot(Deriv(f), a, b, n);  // plot f' over [a,b]
  plot(Integral(f, x0), a, b, n);
  riemann_sum(f, a, b, n, TYPE);
\end{verbatim}
The second graphs the definite integral \teximath{x\mapsto\int_{x_0}^x
f(t)\,dt} over~\teximath{[a,b]}. As above, \teximath{x_0}~defaults
to~0. The third draws rectangles or trapezoids whose area approximates
the definite integral of~\teximath{f} over~\teximath{[a,b]}. The
\code{TYPE} may be \code{UPPER}, \code{LOWER}, \code{LEFT},
\code{RIGHT}, \code{MIDPT}, or~\code{TRAP}.

Tangent lines and envelopes (families of tangent lines) are drawn with
\findex{tan\_line}
\findex{envelope}
\findex{tan\_field}
\begin{verbatim}
  tan_line(f, t);                // f real- or vector-valued
  envelope(f, t_min, t_max, n);  // family of tangent lines
  tan_field(f, g, t_min, t_max, n); // field of tangents
\end{verbatim}
The sample files \filename{conic.xp}~and \filename{lissajous.xp}
illustrate these features.


\subsubsection*{Systems of Differential Equations}

Let \code{F}~be a \code{P}-valued function of two or three variables.
\findex{ode\_plot}
\findex{flow}
\begin{verbatim}
  ode_plot(F, p_0, t_min, t_max, n); 
  flow(F, p_0, t_max, n);
\end{verbatim}
The first plots the solution curve of the initial-value problem
\teximath{\dot{x}=F(x)}, \teximath{x(0)=p_0}, over the specified time
interval.  If \teximath{t_\mathrm{min}} is omitted, its value is~0, so
the curve starts at~\teximath{p_0}. With manual calculation to rotate
a planar field a quarter turn, \code{ode\_plot} can be used to draw
level curves (isobars) of a function of two variables; see the sample
file \filename{dipole.xp}.  The \code{flow} function returns the
result of starting at~\teximath{p_0} and flowing by~\teximath{F} for
time~\teximath{t_\mathrm{max}}, using Euler's method with~\teximath{n}
time steps.  This is useful for placing markers or arrowheads precisely
along a flow line.

A planar or spatial vector field itself may be plotted over a
domain~\code{R} in three styles:
\findex{vector\_field}
\findex{dart\_field}
\findex{slope\_field}
\begin{verbatim}
  vector_field(F, R, [scale]); //  true length
  dart_field  (F, R, [scale]); // const length
  slope_field (F, R, [scale]); // const length
\end{verbatim}
The field is sampled at the grid points of the coarse mesh. If the
domain is 2-dimensional, the plot is a planar slice of the field, even
if the field depends on three variables. If the domain is
3-dimensional, the field is drawn in successive slices
\teximath{x_3=\mathrm{const}}, starting at the height of the first
corner of~\code{R} and ending at the height of the second corner.

The optional final argument, which defaults to~\teximath{1}, scales the
arrowheads in a vector field, and scales the (constant) length of
field elements for slope and dart fields. The sample files
\filename{layout2.xp}, \filename{lorenz.xp}, \filename{slopefield.xp},
and \filename{vfield.xp} illustrate usage.

In each field-plotting command, the domain argument may be replaced by
two points, representing corners of a coordinate rectangle, and two
integers, the number of grid intervals in the selected coordinate
directions. Only planar slices of a vector field can be plotted using
the alternative syntax.
\index{Plotting|)}
\index{Plotting!calculus|)}


\subsection{Non-Euclidean Geometry}
\label{non-eucl}

\ePiX\ provides limited features for spherical and hyperbolic
geometry: the ability to draw lines in the half-plane and Poincar\'e
disk models of the hyperbolic plane, and to draw latitudes,
longitudes, great circle arcs, spherical triangles, regular polyhedra,
and parametrized curves on a sphere.

Hyperbolic line segments are specified by their endpoints in the upper
half space or ball (Poincar\'e) models.  In each case there is no
output if either endpoint lies outside the model.
\findex{hyperbolic\_line}
\findex{disk\_line}
\begin{verbatim}
  hyperbolic_line(p, q);
  disk_line(p, q);
\end{verbatim}
For compatibility with 2-dimensional hyperbolic space, the half-space
model is the set \teximath{\{(x_1,x_2,x_3)\mid x_2>0\}}.

\index{Plotting!spherical|(}
A \code{frame} determines geographical coordinates on a \code{Sphere}:
the first element points toward longitude~0 on the equator, the third
element points to the north pole.  A latitude line depends on a
\code{Sphere}, a \code{frame}, the numerical latitude, and a range of
longitudes. A longitude line is described similarly.
\findex{latitude}
\findex{longitude}
\findex{back\_latitude}
\findex{back\_longitude}
\begin{verbatim}
  latitude(lat, long_min, long_max, Sphere S, frame coords);
  longitude(lngtd, lat_min, lat_max, Sphere S, frame coords);
\end{verbatim}
By default, \code{coords} is the standard \code{frame} and \code{S} is
the unit sphere. These commands draw only the portion of the curve
that is visible from the current viewpoint. The function
\code{back\_latitude} draws the invisible portion of a latitude line.

Spherical arcs and triangles are described by their endpoints. Only
the direction vector from the center of the sphere to an endpoint is
significant; if a sphere is scaled or moved, the same function call
will draw the corresponding object on the new sphere. 

The following draw the visible (front) portions of great circle arcs:
\findex{front\_arc}
\findex{front\_arc2}
\findex{front\_line}
\findex{back\_arc}
\findex{back\_arc2}
\findex{back\_line}
\begin{verbatim}
  front_arc(p1, p2, S);  // short arc on S from p1 to p2
  front_arc2(p1, p2, S); // arc from p1 to -p1 through p2
  front_line(p1, p2, S); // great circle through p1 and p2
\end{verbatim}
Triangles and regular (Platonic) polyhedra are provided. The sample
file \filename{sample/polyhedra.xp} illustrates usage.
\findex{front\_triangle}
\findex{front\_tetra}
\findex{front\_cube}
\findex{front\_octa}
\findex{front\_dodeca}
\findex{front\_icosa}
\findex{back\_triangle}
\findex{back\_tetra}
\findex{back\_cube}
\findex{back\_octa}
\findex{back\_dodeca}
\findex{back\_icosa}
\begin{verbatim}
  front_triangle(p1, p2, p3, S); // spherical triangle
  front_tetra(S, coords);  // regular tetrahedron
  front_cube(S, coords);   //    hexahedron
  front_octa(S, coords);   //    octahedron
  front_dodeca(S, coords); //    dodecahedron
  front_icosa(S, coords);  //    icosahedron
\end{verbatim}
Each function has a \code{back} version, which draws the hidden
portion. The tetrahedron, cube, and octahedron are (up to scale)
inscribed in the cube of side length~\teximath{2} centered at the origin whose
sides are parallel to~\code{frame}. The point \teximath{(1,1,1)} is a vertex of
the tetrahedron.

Up to scale, the icosahedron's vertices lie on the golden rectangle
with vertices \teximath{(\pm\gamma,0,\pm1)} and its images under cyclic
permutation of coordinates. The dodecahedron is dual to the
icosahedron.

\subsubsection*{Spherical Geometry}

\index{Sline}

The \code{Sline} class represents geodesic arcs on the unit sphere,
and provides intersection and reflection operators.
\findex{Sline}
\findex{pole}
\findex{cosine}
\findex{reflect}
\findex{collinear}
\findex{cuts}
\findex{draw}
\findex{draw\_front}
\findex{draw\_back}
\findex{draw\_line}
\findex{line\_front}
\findex{line\_back}
\begin{verbatim}
  Sline L1(tail, head); // non-proportional vectors
  L1.pole();            // pole of equator
  L1.cosine();          // subtended arc
  L1.reflect(arg);      // reflect arg across L1
  L1.reflect(L2);       // reflect L2 across L1
  L1.collinear(L2);     // test for collinearity
  L1.cuts(L2);          //   or crossing
  L1.draw();            // arc
  L1.draw_front();      // front arc (same for back)
  L1.draw_line();       // draw entire line
  L1.line_front();      // front line (same for back)
\end{verbatim}


\subsubsection*{Spherical Plotting}

Parametrized paths on a \code{Sphere~S} (by default the unit sphere)
can be specified either by radial projection of a \code{P}-valued
curve, or by stereographic projection of a plane curve given as a pair
of \code{double}-valued functions:
\findex{frontplot\_R}
\findex{frontplot\_N}
\findex{frontplot\_S}
\findex{backplot\_R}
\findex{backplot\_N}
\findex{backplot\_S}
\begin{verbatim}
  frontplot_R(phi, t_min, t_max, n, [S]);    // radial
  frontplot_N(f1, f2, t_min, t_max, n, [S]); // from north pole
  frontplot_S(f1, f2, t_min, t_max, n, [S]); // from south pole
\end{verbatim}
Attempts to perform radial projection on a path through the origin
will generate division-by-zero errors. Stereographic projection maps
the equatorial plane \teximath{\{x_3=0\}} to the unit sphere by
projection from the corresponding pole: \teximath{N=(0,0,1)},
\teximath{S=(0,0,-1)}.

Each spherical plot command has a \code{back} version that prints only
the portion of the path invisible from the current viewpoint.
Because of the way \ePiX\ layers output, it is generally best to put
hidden portions of the input before visible portions, with line width
and/or style that suggests hidden lines.
\index{Plotting!spherical|)}


\subsection{Data Plotting} 
\index{Plotting!data|(}
\index{Data plotting|(}

Thanks in large part to code and ideas from Marcus Hanwell, files of
numerical data can be created, manipulated, analyzed, plotted (paths,
scatter plots, and histograms), read, and written.  The format for a
data file is one or more floating-point numbers per line, with the
same number of entries per line. Anything that appears on a line after
the \LaTeX\ comment character~\code{\%} is a comment.

\ePiX\ provides two \code{plot} commands for file data. The first
facilitates plotting selected columns; the second simplifies plotting
the first two columns with one or both scales logarithmic. Either form
can be used to plot selected columns logarithmically.

The general commands read numbers from two or three columns of a
specified file, pass them as arguments to a \code{P}-valued
function~\code{F}, and plot the resulting points:
\findex{plot}
\begin{verbatim}
  plot("filename", STYLE, [i_1], [i_2], [i_3], [F]);
  plot("filename", STYLE, F, [i_1], [i_2], [i_3]);
\end{verbatim}
The first argument is the name of the data file. The \code{STYLE} may
be \code{PATH}, which joins the points in the order they appear, or
any of the marker types in Table~\ref{table-marker}.  The
integers~\teximath{i_k} specify columns from which to extract data; these
default, respectively, to the first column, second column, and null (a
column of zeroes). If the ``coordinate system''~\code{F} is omitted in
the first command, it defaults to the Cartesian point constructor. The
function~\code{F} is mandatory in the second form; useful choices
include \code{log\_log}, \code{log\_lin}, and~\code{lin\_log}, which
plot the corresponding coordinate logarithmically.


\subsubsection*{Data Files}

For more elaborate analysis, the \code{data\_file} class presents an
interface to a file as an ordered list of columns. There are two
general ways to create a \code{data\_file}: read in an external file,
or generate data (up to three columns) using specified
\code{double}-valued functions. In the constructors below, each
function~\code{fi} is a \code{double}-valued function of one variable.
\findex{data\_file}
\findex{data\_file.read}
\begin{verbatim}
  data_file DF("my_data"); // read data from disk file
  data_file DF(f1, t_min, t_max, num_pts); // values of f1
  data_file DF(f1, f2, t_min, t_max, num_pts);
  data_file DF(f1, f2, f3, t_min, t_max, num_pts);
  data_file DF(3);  // create empty data_file with 3 columns
  DF.read("file1"); // read a disk file
\end{verbatim}

Columns of a \code{data\_file} can be transformed by a user-specified
function, averaged, correlated, extracted (for use by other code),
scatter plotted, and written to a disk file at specified
precision. Below, the function~\code{f} is a \code{double}-valued
function of one variable and~\code{F} is a \code{P}-valued function of
two or three variables, whose components are written back to the
selected columns.
\findex{data\_file.transform}
\begin{verbatim}
  DF.transform(f, i);    // apply f to selected column(s)
  DF.transform(F, i=1, j=2);
  DF.transform(F, i, j, k);
\end{verbatim}
Basic statistical operations on columns are provided.
\findex{data\_file.dot}
\findex{data\_file.avg}
\findex{data\_file.var}
\findex{data\_file.covar}
\findex{data\_file.regression}
\begin{verbatim}
  DF.dot(i,j);        // dot product of columns i, j
  DF.avg(i);          // mean of column i
  DF.var(i);          // population variance
  DF.covar(i,j);      // covariance
  DF.regression(i,j); // plot regression line
\end{verbatim}
A \code{data\_file} is scatter plotted using syntax as described above.
% The output is cropped automatically to the bounding box.
\findex{data\_file.plot}
\begin{verbatim}
  DF.plot(STYLE, [i1], [i2], [i3], [F]);
  DF.plot(STYLE, F, [i1], [i2], [i3]);
\end{verbatim}
Histograms and bar charts are described below.

A \code{data\_file} can be written to a disk file as raw data, or in
specified format. Below, \code{fmt} denotes a \code{string}-valued
formatting function of two variables and \filename{myfile} is the
name of the disk file to be written.
\findex{data\_file.precision}
\findex{data\_file.write}
\begin{verbatim}
  DF.precision(4);    // set to 4 significant figures
  DF.write("myfile"); // write as tab-separated columns
  DF.write("myfile", fmt, [i1], [i2]); // apply fmt to cols
\end{verbatim}
A column can be extracted as a \code{C++}~\code{vector} for use by
another function.
\findex{data\_file.column}
\begin{verbatim}
  DF.column(i);    // i-th column
  DF.column(f, i); // i-th column, transformed by f
\end{verbatim}


\subsubsection*{Data Containers}

\ePiX\ provides a \code{data\_mask} class for culling data from a file
according to the values in a specified column, and a \code{data\_bins}
class for sorting data by value.

A \code{data\_mask} consists of an interval of numbers and a
``filter'' function. The (closed, open, or half-open) interval is
given as a string in standard mathematical notation, or by its
endpoints (for a closed interval).  The filter is a
\code{double}-valued function of \code{double}, by default the
identity, \teximath{f(x)=x}.
\findex{data\_mask}
\begin{verbatim}
  data_mask dm("[a, b]", [f]);
  data_mask dm(a, b, [f]);
\end{verbatim}
A \code{data\_mask} ``passes'' inputs~\teximath{x} if~\teximath{f(x)}
lies in the interval. The member function \code{reverse()} inverts
this logical test.

The \code{data\_file} class has \code{prune} functions to cull
rows for which a specified column's entry satisfies a
\code{data\_mask}'s criterion.
\findex{data\_mask.prune}
\begin{verbatim}
  DF.prune(dm, i);   // remove row if i-th column entry fails
  DF.prune(a, b, i); // remove row if i-th column outside [a,b]
\end{verbatim}

A \code{data\_bins} object models an interval divided at specified
locations into ``bins'', not necessarily of equal length. Numerical
data is read in and the number of points in each bin counted. The
lifetime of a \code{data\_bins} object has two stages. First, ``cuts''
(endpoints of subintervals) are added. Then, once data is read, the
cuts are ``locked'' and cannot subsequently be changed.
\findex{data\_bins}
\findex{data\_bins.cut}
\findex{data\_bins.read}
\begin{verbatim}
  // [xmin, xmax] divided into n equal intervals, 1 by default
  data_bins db(xmin, xmax, [n]);
  db.cut(x);  // add a cut at x (if x is in bounds)
  db.read(vector<double>); // read data, lock bins
\end{verbatim}

\index{Histograms}
\index{Bar Charts}
A \code{data\_bins} object can be plotted as a histogram (rectangles'
\emph{area} is proportional to the bin population), bar chart
(rectangles' \emph{height} is proportional to the bin population), or
spline interpolation of a bar chart.
\findex{data\_bins.histogram}
\findex{data\_bins.bar\_chart}
\findex{data\_bins.plot}
\begin{verbatim}
  db.histogram(c=1);  // c = vertical scale factor
  db.bar_chart(c=1);
  db.plot(c=1);
\end{verbatim}
By default (\code{c=1}), the height of a bar chart rectangle is the
fraction of the total population contained in the bin; thus, the
height is always between 0~and 1. For a histogram, the height of a
rectangle is the fraction of the population per horizontal unit in the
bin; thus, the total area over an interval~\teximath{[a,b]} does not
depend on how~\teximath{[a,b]} has been subdivided. The sample files
\filename{dataplot.xp}~and \filename{histogram.xp} illustrate use.

Statistical convention dictates cuts be chosen distinctly from data
values; that is, values should all fall strictly within a bin, not at
a boundary point. With large, unknown data sets, this convention may
be difficult to uphold. \ePiX\ attempts to handle anomalous data
intelligently, by keeping counts of values ``out-of-range'' or
``on-cut''.

In detail, if \teximath{x<a} or \teximath{x>b} is a data value, it is
counted as out-of-range and does not contribute to the histogram
population. If \teximath{x=a} or \teximath{x=b}, the value counts as
both out-of-range and on-cut but is added to the population of the
lowest or highest bin, respectively. Any other cut appearing as a data
input is flagged as on-cut, and increments the population of each
adjacent bin by one-half. When a histogram or bar chart is written,
\ePiX\ prints a warning message summarizing the number of anomalous
data seen.


\subsubsection*{Error Bars}
\index{Error bars}

Simple horizontal and vertical error bars are provided.  The final
(optional) argument is the true height or width (respectively) in pt.
\findex{h\_error\_bar}
\findex{v\_error\_bar}
\begin{verbatim}
  h_error_bar(P location, double error, <mark type>, ht=6);
  v_error_bar(P location, double error, <mark type>, wd=6);
\end{verbatim}

To create more complex elements, such as asymmetric bars, whisker
plots, labeled error bars, and the like, write a custom function using
true-size drawing, see page~\pageref{truesize}. For example, a
fillable, labeled, asymmetric, vertical rectangular error bar can be
implemented (entirely in \ePiX) like this:
\begin{verbatim}
  void error_bar(P loc, double lo, double hi, P offset,
                 std::string text, align, double wd=6)
  {
    const double width(pt_to_screen(0.5*wd)); // converts to 3pt
    rect(loc - P(width, lo), loc + P(width, hi));
    line(loc - P(width,  0), loc + P(width,  0));
    label(loc + P(width, 0), offset, text, align);
  }
\end{verbatim}
For stylistic uniformity, functions such as this should be put into a
library and used systematically. Section~\ref{section-extensions}
outlines the process of writing, compiling, and using a custom
library.
\index{Plotting!data|)}
\index{Data plotting|)}


\subsection{Legends}
\index{Legends}

A \code{legend} systematically labels different parts of a plot by
associating visual ``keys'' with explanatory text. This tends to be
most useful for plots containing several distinct but conceptually
related graphs requiring contrast. 

Visually, a \code{legend} is represented as an aligned list of rows,
each containing a box (the key), a gap (the label skip), and some
text.  These rows are printed in a (usually large) masked label. By
default, keys are \code{12pt}~squares bordered in black, the label
skip is~\code{6pt}, the background is white, and the border is plain
black. These attributes are controlled (simultaneously for all items)
with member functions.  Parameters of type \code{double} represent
lengths in~\code{pt}.
\findex{legend}
\findex{legend.backing}
\findex{legend.border}
\findex{legend.item\_border}
\findex{legend.label\_skip}
\findex{legend.key\_size}
\begin{verbatim}
  legend L;
  L.backing(color);          // set background
  L.border(color, [double]); // set border color [and width]
  L.border(double);          // border width

  L.item_border(color, [double]);
  L.item_border(double);     // same, for item borders

  L.label_skip(double);
  L.key_size(double);
\end{verbatim}


\subsubsection*{Legend Keys}

There are three types of legend key, representing filled regions,
paths, and markers. Fill and path keys get their visual attributes
from the current drawing state. A mark key must be told the marker
type. Each is created by specifying the item text.
\findex{legend.fill\_item}
\findex{legend.path\_item}
\findex{legend.mark\_item}
\begin{verbatim}
  L.fill_item(text);
  L.path_item(text);
  L.mark_item(<mark type>, text);
\end{verbatim}
Keys in the printed legend appear in the same order they are created
in the input file.


\subsubsection*{Creating a Legend}

Normally, a \code{legend} is defined near the start of a file, and an
item is added at the point in the file where the corresponding object
or plot is drawn, so that the item receives the correct attributes. A
\code{legend} is placed into the figure with the \code{draw}
function. The arguments have the same meaning as for ordinary
\code{label}s.
\findex{legend.draw}
\begin{verbatim}
  L.draw(P location, P offset, align);
\end{verbatim}
All three arguments are mandatory.  ``Global'' \code{legend} settings
(border, backing, etc.) may be changed anywhere between the
\code{legend}'s creation and \code{draw} function. The sample file
\filename{legend.xp} revisits the example on
page~\pageref{taylor_poly}, including a trig-labeled axis and a
legend. The file \filename{shadeplot.xp} contains filled keys, and
\filename{dataplot.xp} contains a \code{legend} with items of mixed
type.


\section{More About \CXX}
\index{Cplusplus@\CXX|(emph}

A textbook or similarly detailed reference is essential for serious
study of~\code{C} or~\CXX. \emph{The \code{C} Programming Language},
second edition, by Kernighan~and Ritchie~\cite{KnR}, is an excellent,
manageable resource for the basics of procedural
programming. \emph{C++ Primer Plus}, by Stephen Prata~\cite{Prata},
clearly lays out the extensive details of~\CXX.  Marshall Cline's
\emph{\CXX\ FAQ Lite}~\cite{Cline} engagingly discusses common points
of confusion and furnishes tips on good design and programming.

\CXX\ is a powerful, complex language whose syntax is similar to that
of~\C, or to the scripting languages of Maple~and Mathematica. An
\ePiX\ input file is source code for a \CXX~program that writes a
\LaTeX\ picture as output.  \ePiX\ may be viewed as an extension
to~\CXX; in the same way that \LaTeX\ furnishes a high-level interface
to~\TeX, \ePiX\ provides a high-level bridge between the computational
power of~\CXX\ and the \LaTeX\ \code{picture} environment.

\index{Function}
\index{Conditional statement}
Like all high-level programming languages, \CXX\ provides variables,
functions, and control structures. Variables hold pieces of data such
as numerical values and geometric locations, while functions operate
on data. A control structure, such as a loop or conditional statement,
affects the program's course according to the program's current state.
A source file is composed primarily of ``statements'', which perform
actions ranging from defining variables and functions to setting
figure attributes, performing calculations, and writing objects to the
output file.

\subsection{Names and Types}

\index{Variable!names|(}
Names of variables and functions may consist (only) of letters,
digits, and the underscore character. The first character of a name
must not be a digit, and the language standard reserves names starting
with underscore for library authors. Names are case-sensitive, but
it's usually a bad idea to use a single name capitalized and
uncapitalized in a single file. Numerous capitalization conventions
are used informally; this document uses uncapitalized words separated
by underscores for variables and functions, and occasionally uses all
capitals for constants. As with names of \LaTeX\ macros, primary
considerations are clarity (of meaning), readability, and consistency.
\index{Variable!names|)}

\index{Variable!type of}
Every variable in \CXX\ has a ``type'',
such as integer (\code{int}), double-precision floating point
(\code{double}), or Boolean (\code{bool}, true~or false).  \ePiX\ 
provides additional types, the most common of which is~\code{P}, for
point. The construct \code{P(x,y,z)} creates~\teximath{(x,y,z)}, while
\code{P(x,y)} gives~\teximath{(x,y,0)}, which is effectively the
pair~\teximath{(x,y)}. A variable is defined by giving its type, its
name, and an initializing expression.

\index{Variable!pointer}
\index{Pointer}
In~\code{C} and~\CXX, a \emph{pointer} variable holds the memory
address of another variable. Pointers are useful for manipulating
(possibly large) data structures through ``handles'' of fixed small
size. \CXX\ also provides \emph{references}, which bind an additional
name to an existing object and allow the object to be manipulated
through this alias. The statements
\begin{verbatim}
  double  x=1;   // ordinary variable definition
  double& rx=x;  // bind a reference, signified by &
\end{verbatim}
define a variable~\code{x} having the value~1, and bind a reference
variable~\code{rx} to it. As long as~\code{rx} exists, it refers
to~\code{x}. If the value of~\code{x} changes, the value of~\code{rx}
does as well. Conversely, the value of~\code{x} can be altered by
assigning to~\code{rx}. However,~\code{rx} is the size of a pointer,
regardless of the size of~\code{x}, so~\code{rx} can be passed
efficiently in a function call.
\texicomment
Some applications are discussed on page~\pageref{references}.
\texicommentend


\subsection{Functions}

\index{Function|(}
In a programming language, the term ``function'' refers to a block of
code that is executable by name.  A \CXX\ function takes a list of
``arguments'', and has a ``return value''.  This information, together
with the function's name, must be provided when a function is defined.
A function may not be defined inside another function. However, a
function may call other functions (including itself) as part of its
execution:
\begin{verbatim}
  int factorial(unsigned int n)
  {
    if (n == 0) return 1;
    else return n*factorial(n-1);
  }
\end{verbatim}

\index{Function!returning \code{void}}
The special type~\code{void} represents a ``null type''. A function
that performs an action but does not return a value has return type
\code{void}. A function that takes no arguments may be viewed as
taking a single \code{void} argument.

Every \CXX\ program has a special function~\code{main()}, which is
called by the operating system when the program is run. The arguments
of~\code{main()} are command-line arguments, and the return type is an
integer that signals success or failure. User-specified functions must
be defined before the call to \code{main()} or in a
separately-compiled file.

Functions in~\CXX\ may be as simple as an algebraic formula or as
complex as an arbitrary algorithm. Greatest common divisors, finite
sums, numerical derivatives and integrals, solutions of differential
equations, recursively generated fractal curves, and curves of best
fit are a few applications in \ePiX. Several sample files contain
user-level algorithms, which do not require knowledge of \ePiX's
internal data structures. The source file \filename{functions.cc}
contains simple functions defined by algorithms, and
\filename{functions.h} illustrates the use of \CXX\ templates. Other
source files, such as \filename{plots.cc}, may be consulted for
Simpson's rule, Euler's method, and the like.


\subsection{Mathematical Functions}

\index{Function!mathematical|(}
\CXX~knows several familiar mathematical functions by name:
\begin{verbatim}
  sqrt   exp   log   log10   ceil   floor   fabs
\end{verbatim}
(\code{fabs} is the absolute value for a floating-point argument.)
\ePiX\ provides trig and inverse trig functions sensitive to angular
mode:
\index{Angle!mode}
\findex{Cos}
\findex{Sin}
\findex{Tan}
\findex{Sec}
\findex{Csc}
\findex{Cot}
\findex{Acos}
\findex{Asin}
\findex{Atan}
\begin{verbatim}
  Cos       Sin       Tan
  Sec       Csc       Cot
  Acos      Asin      Atan
\end{verbatim}
The inverse functions are principle branches.

The function \code{pow(x,y)} returns~\teximath{x^y} when
\teximath{x>0}, and \code{atan2(y,x)} (N.B. argument order) returns
\teximath{\mathrm{Arg}(x+iy)\in(-\pi,\pi]}, the principle branch of
arg.  \CXX~knows many constants to 20~decimal places, such as
\code{M\_PI}, \code{M\_PI\_2}, and \code{M\_E} for~\teximath{\pi},
\teximath{\pi/2}, and~\teximath{e} respectively.  \ePiX\ defines a few
additional functions:
\findex{sgn}
\findex{zero}
\findex{sinx}
\findex{cb}
\begin{verbatim}
  sgn   zero   sinx   cb
\end{verbatim}
\texicomment
\begin{center}
  \input{sgn.eepic}\qquad
  \input{sinx.eepic}\qquad
  \input{cb.eepic}
\end{center}
\medskip
\texicommentend
\code{sgn} is the signum function; \code{zero} is the constant
function; \code{sinx} is the function \teximath{x\mapsto \sin(x)/x}
with the discontinuity removed; \code{cb} (for ``Charlie Brown'') is
the period-2 extension of the absolute value function
on~\teximath{[-1,1]}.

The GNU \CXX\ library defines other functions, including inverse
hyperbolic functions (\code{acosh}, etc.), \code{log}~and \code{exp}
with base~2, 10, or arbitrary~\teximath{b} (\code{log2}, etc.), the
error and gamma functions (\code{erf}~and \code{tgamma} [sic],
respectively), and Bessel functions of first and second kind:
\code{j0}, \code{j1}, \code{y0}, etc. Use, e.g., \code{jn(5,~)} to get
higher indices.  The GNU \C~library reference manual~\cite{GNUC}
describes these and other functions in detail.

Functions may be used in subsequent definitions, and functions of two
(or more) variables are defined in direct analogy to functions of one
variable:
\begin{verbatim}
double f(double t) { return t*t*log(t*t); } // t^2 \ln(t^2)
double g(double s, double t) { return exp(2*s)*Sin(t); }
\end{verbatim}
\index{Function!mathematical|)}


\subsection{Basics of Classes}

\index{Classes|(}
Unlike~\code{C}, \CXX\ supports ``object-oriented programming''. In a
nutshell, a \emph{class} is an abstraction in computer code of some
concept, such as a point, a sphere, a mapping that can be plotted, or
a camera. Classes allow a programmer to separate an object's
\emph{interface} (the set of meaningful operations) from its
\emph{implementation} (the data structures and algorithms that realize
the interface). 

\index{Function!class member}
A class implementation consists of \emph{members} (named data
elements) and \emph{member functions} (functions that belong to the
class and have free access to members).  \CXX\ classes enforce access
permissions on their members, protecting data from being manipulated
except as promised by the interface.

An ideal interface looks like a black box: It hides the implementation
completely.  In order to cooperate, two classes need only know each
other's interfaces. This separation of form and function modularizes a
program, and facilitates debugging, code reuse, and overall
maintainability, particularly in large programs. 

In simple programming, classes may be treated like built-in types.
Each class object has its own member functions, whose call syntax
differs from standard function calls:
\begin{verbatim}
  Circle C1(P(1,0), 1.5); // circle of given center and radius
  C1.draw();              // member function Circle::draw();
\end{verbatim}
Naturally, this call draws the circle~\code{C1}. Generally, a member
function call consists of a class object's name, a period, and the
name of the member function. Arguments, if any, go in the parentheses
after the member function name, just as in a regular function call.

A few short paragraphs cannot do more than scratch the surface of
classes and object-oriented programming. For more details, please
consult a book, such as Prata~\cite{Prata} or Stroustrup~\cite{S}, or
Cline's on-line FAQ~\cite{Cline}.
\index{Classes|)}


\subsection{References and Function Arguments}

\index{Function!call syntax}
\index{Variable!reference}
\label{references}
\code{C}~and \CXX\ are ``call by value'' languages. Variables are not
passed to a function; instead a copy of the value is made, and the
function operates only on the copy. Though this feature causes
occasional inconvenience, it prevents an object from being altered
unexpectedly by a function call in a different part of the program.
Calling by value helps localize the logic of a program, and
circumvents easy-to-write but extremely hard-to-find bugs.

In \CXX, a function may accept reference arguments. Passing an object
by reference grants the calling function access to the object itself,
not to a copy. There are two common applications: The object is a
large data structure for which copying is ``expensive'', or the
function \emph{needs} to modify its arguments (e.g., a function
\code{swap(x,y)} that exchanges the values of \code{x}~and
\code{y}). 

For the first situation, \CXX\ provides the \code{const} keyword,
which ensures the function does not modify its arguments, but accepts
a reference merely for efficiency. Any attempt to modify a
\code{const} argument will be caught by the compiler. Most \ePiX\
commands accept \code{const} reference arguments.

The ability to pass function arguments by reference is sometimes
touted as a feature in \CXX\ texts. However, the technique circumvents
the data encapsulation of calling by value, and should be avoided
unless absolutely necessary. If a function merely ``updates'' the
value of a variable, probably the variable should be of class type,
and the update should be performed by a member function.

A function declaration must indicate that its arguments are
references. The declarations below have the indicated idiomatic
meanings.
\begin{verbatim}
  class matrix;
  double det(matrix); // call by value, perhaps inefficient
  matrix& transpose(matrix&); // probably changes its argument
  double trace(const matrix&); // does not change its argument
\end{verbatim}
Unlike pointer arguments, reference arguments impose no syntactic
burden on the user. If~\code{A} is a~\code{matrix}, then
\code{transpose(A);} and \code{trace(A);} will compile. You need not
declare explicit reference variables and pass those to the function.


\subsection{Overloading}

\index{Function!overloaded}
\CXX\ provides ``overloading'': Multiple functions can be given the
same name, so long as the number and/or type of their arguments
differ. (It is \emph{not} enough for the return types alone to differ.
The compiler must be able to select a function from its calling
syntax.)  To the user, the appearance is that a single function
intelligently handles multiple argument lists. Naturally, overloaded
names should refer to functions that are conceptually related.
\index{Function|)}


\subsection{Scope}

\index{Variable!local}
A \CXX\ statement ends with a semicolon. A collection of statements
enclosed by curly braces is a ``code block'', and may be viewed as a
single logical statement. Curly braces determine a ``scope'', inside
which variable names may be re-used without ambiguity.  Function
bodies are code blocks, as are the alternatives associated to control
statements.  A variable defined between curly braces is said to be
\emph{local} to the scope in which it is defined; its value cannot
be used out of scope.  Variables should be declared in the smallest
scope possible.

The compiler is not picky about spaces, tabs, and newlines, so an
input file should be formatted to make local scopes visually
apparent. \code{emacs} automatically indents code to reflect scope,
though the default behavior does not please all users.  As with
variable naming, clarity and consistency are paramount.


\subsection{Headers and Pre-Processing}

A \CXX\ source file is compiled in multiple stages that occur
transparently to the user. The first step, pre-processing, involves
simple text replacement for file inclusion, macro expansion and
conditional compilation. Next, the source is compiled and assembled:
Human-readable language instructions are parsed, then represented in
assembly language. Finally, the object files are linked: Function
calls are resolved to hard-coded file offsets, possibly involving
external library files, and the program instructions are packaged into
an executable binary that the operating system can run.

Pre-processing is used much less in~\CXX\ than in~\code{C}; the
language itself supports safer and more featureful alternatives to
macros, such as \code{const} variables and inline functions. File
inclusion and conditional compilation are the chief uses of the
pre-processor. Lines of the form 
\begin{verbatim}
  #include <iostream>
  #include "epix.h"
\end{verbatim}
cause the contents of a \emph{header file} to be read into the source
file. A header file contains variable and function
\emph{declarations}, statements that specify types and names but do
not define actual data. Declarations tell the compiler just enough
to resolve expressions and function calls without knowing specific
values or function definitions. 

\index{Input file!conditional compilation}
Conditional compilation is similar to conditional \LaTeX\ code. For
example, a file might produce either color or monochrome output as
follows:
\begin{verbatim}
#ifdef COLOR
  ...  // code for generating color figure
#endif /* COLOR */
#ifndef COLOR
  ...  // monochrome code
#endif /* undef COLOR */
\end{verbatim}
The ``compiler symbol'' \code{COLOR} is an ordinary \CXX\ name. To
control compilation, either put a \code{\#define COLOR} line in the
file, or (better) supply the flag on the command line: \\
\verb+  epix -DCOLOR <file.xp>+ \\
Every \code{\#ifdef} must have a matching \code{\#endif}. Commenting
the \code{\#endif} is a good habit; in a realistic file, the start and
end of a conditional block may be separated by more than one screen.


\subsection{Comparison with \LaTeX\ Syntax}

As a programming language, \CXX\ provides certain features common to
all languages (such as \LaTeX, MetaPost, Perl, Lisp\ldots) and adheres
to rules of grammar. Salient differences between \LaTeX\ and \CXX\ 
include:

\begin{enumerate}
\item Every \CXX\ statement and function call must end with a
  semicolon. An omitted semicolon may result in a cryptic error
  message from the compiler. Pre-processor directives, which start
  with a~\code{\#}, do not end with a semicolon.

\item Backslash is an escape character in \CXX:
  \index{Label}
\begin{verbatim}
  // Put label $y=\sin x$ at (2,1)
  // Note single  ^ backslash in output
  label(P(2,1), P(0,0), "$y=\\sin x$"); 
  //       Double backslash ^^ in source
\end{verbatim}
  
\item Variable and function names may contain letters (including
  underscore) and digits \emph{only}, are case sensitive, and must
  begin with a letter.

\item Variables in \CXX\ must have a declared \emph{type}, such as
  \code{int} (integer) or \code{double} (double-precision floating
  point). If a variable has global scope and its value does not
  change, the definition should probably come in the preamble or at
  the beginning of \code{main}. Local variables should be defined in
  the smallest possible scope. Unlike~\C, \CXX\ allows variables to be
  defined where they first appear.
  
\item \CXX~requires explicit use of~\code{*} to denote multiplication;
  juxtaposition is not enough. \CXX~does not support the use of
  \verb+^+ for exponentiation, e.g., \verb+t^2+ is invalid. Instead,
  use \code{t*t}~or \code{pow(t,2)}.

\item \CXX\ has single- and multi-line comments. Everything between a
  double slash and the next newline is ignored, while the strings
  \verb+/*+~and \verb+*/+ delimit multi-line comments. A single-line
  comment may appear within a multi-line comment, but the compiler
  does not nest multi-line comments.
  \index{Input file!comment in}

\end{enumerate}

Between them, \C~and \CXX~have about 100~reserved keywords which
cannot be used as function or variable names.
\index{Cplusplus@\CXX|)}


\section{Attribute Quick Reference}
\label{section-quickref}

In the body of an input file, the ``drawing state'' determines the
figure's appearance. Attributes are declarations, set by commands that
accept arguments of the stated type.

A \code{len} argument is a double-quoted string containing a number
and a two-letter \LaTeX\ length unit, such as \code{"1.5pt"} or
\code{"6cm"}. A \code{color} argument is a named primary
(\code{Red()}, \code{Cyan()}, \code{White()}, etc.), a \code{Color}
specified by densities (\code{RGB(...)}, \code{CMYK(...)}, etc.), or a
\code{Color} object. Using \code{Neutral()} as a \code{Color} argument
generally turns off the corresponding attribute.

\begin{itemize}

\item Angular mode: \code{radians()}, \code{degrees()}, or
  \code{revolutions()}.

  The angular mode affects all trigonometric operations, including
  camera rotations, the drawing of arcs and ellipses, polar plotting,
  label angle, and the trig functions themselves. Angle-sensitive trig
  functions are capitalized, e.g., \code{Cos}, \code{Tan}.

\item Fill style: \code{fill(color)}, \code{fill(bool)},
  \code{nofill()}.

\item Path style:

  \begin{itemize}

  \item Width: \code{plain()}, \code{bold()}, \code{bbold()},
    \code{pen(len)}.

  \item Line style: \code{line\_style(string)}. The argument is a
    WYSIWYG sequence of dashes, spaces, and periods.
    \code{dash\_size(double)} and \code{dot\_sep(double)} set the
    (approximate) length in~\code{pt} of the pattern.  The commands
    \code{solid()}, \code{dashed()}, \code{dotted()} define
    ``standard'' defaults for brevity.

  \item Path color:
    Paths can be drawn using \emph{two} pens, one atop the other. When
    the ``base'' pen is white (or the background color) and wider than
    the ``line'' pen, a path masks parts of the figure it crosses. A
    3-D effect may be obtained by making the base pen a darker shade
    than the ``line'' pen.

    \code{pen(color,[len])}, \code{base(color,[len])}

    Standard widths: \code{plain(color)}, \code{bold(color)},
    \code{bbold(color)}

  \end{itemize}

\item Text attributes:

  \begin{itemize}

  \item Color: \code{label\_color(color)}

  \item Mask: \code{label\_mask(color)}, \code{label\_pad(len)}

  \item Border: \code{label\_border(color, [len])},
    \code{label\_border(len)}. The command \code{no\_label\_border()}
    turns off label borders.

  \item Font size: \code{font\_size(LaTeX size)}, no argument means
    \code{normalsize}.

  \item Font face: \code{font\_face(LaTeX font)}, two-letter font
    selection string, default is \code{rm}.

  \item Rotation: \code{label\_angle(double)}

  \end{itemize}

\end{itemize}

Do not confuse \code{Color} constructors with the similarly-named
(deprecated) lowercase attribute-setting commands, \code{rgb(r,g,b)},
\code{cmyk(c,m,y,k)}, \code{red(d)}, etc. These commands affect text,
paths, and filled regions. For example, the single command
\code{red()} has the same effect as the three commands
\code{pen(Red())}, \code{fill(Red())}, and \code{label\_color(Red())}.


\texinfochapter{4}{Advanced Topics}
\label{chapter-adv}

This chapter covers \emph{ad hoc} tricks and open-ended techniques
that require relatively more programming sophistication. You will
almost surely need an external \CXX\ reference if you do not speak the
language.


\section{Hidden Object Removal}
\label{section-hidden}

\index{Hidden object removal|(}
\ePiX\ writes the output file in the same order that objects appear in
the input. The order is significant because PostScript builds a figure
in layers: Objects are drawn over objects that come earlier in the
file. Shaded polygons can be used to obtain surprisingly effective
hidden object removal in surface meshes. This section describes the
data structures defined in the source files \code{surface.*}.

\index{Camera}
The basic idea is to create a shaded polygon class that knows its
approximate distance to the camera. For computational simplicity, a
mesh ``facet'' is treated as a quadrilateral, located at the
arithmetic mean of its vertices. A facet's boundary is created from a
map and a domain by tracing a fine mesh rectangle counterclockwise.

To draw a parametrized surface, facets are stored in a \CXX\ vector,
sorted in decreasing order of distance to the camera, and printed to
the output file. If filling is active, the gray density of a facet
depends on the cosine of the angle between the normal vector and the
vector from the camera to the element.

This simple algorithm works surprisingly well when mesh elements
intersect at most along complete edges. To incorporate line-like
elements (e.g., coordinate axes, wire-mesh plots) with shaded
surfaces, the best technique is often to order high-level scene
elements manually, breaking up shaded surfaces (for example, with
domain resizing or clipping) as necessary. The sample file
\filename{saddle.xp} illustrates possible techniques.

Shaded surfaces can be decorated with a bit of hackery. For example,
the \code{facet::draw} function in \filename{facet.cc} can be
modified easily to draw line elements, tangents, or normal vectors
along with the facet itself. The sample file
\filename{decorate.xp} contains a couple of ideas. (The
decorations are activated by compiler flags; please consult the file
itself for information on compiling.)
\index{Hidden object removal|)}


\section{Extensions}
\label{section-extensions}

Thanks to a suggestion of Andrew Sterian, \ePiX\ is extensible. User
extensions span a spectrum, from header files that require only basic
knowledge of~\CXX\ to separately compiled libraries that add
substantial new features.  The structure of the source code is
outlined in Section~\ref{section-programmers}.


\subsection{Header Files}

A \CXX~header file conventionally has suffix~\filename{.h}, as in
\filename{myheader.h}. To use this custom header, put a line
\code{\#include "myheader.h"} in your source file.

User definitions can be easily and robustly implemented with ``inline
functions''. Inline functions are superficially similar to macros, but
are far more safe and featureful (since they are handled by the
compiler rather than by the pre-processor).  Examples are
\begin{verbatim}
  inline void Bold() { pen(1.6); }
  inline void purple() { rgb(0.5, 0, 0.7); }
  inline void draw_square(double s) { rect(P(-s,-s),P(s,s)); }
  inline double cube(double x) { return pow(x,3); // x^3 }
\end{verbatim}
The keyword \code{void} signifies a function that does not return a
value, or (when used as an implicit parameter) a function that does
not accept arguments. Inline function definitions are syntactically
identical to ordinary function definitions, but \emph{must} occur in a
header file or in the source file where they are used. The examples
above might be used in an input file as follows:
\begin{verbatim}
  Bold();  
  draw_square(cube(1.25));
\end{verbatim}


\subsection{Compiling}

The next few sections outline the creation of a ``static library'' on
GNU/Linux, and explain how to incorporate custom features at runtime.
The extensively-commented sample files \filename{std\_F.cc}~and
\filename{std\_F.h} illustrate the techniques described below, and may
be used for guidance and experimentation.

A small library is usually written as a \emph{header} file, which
contains class and function declarations (also called ``prototypes''),
and a \emph{source} file, which contains the actual
code. Conventionally (under *nix), these files have extension
\filename{.h}~and \filename{.cc} respectively. Header and source files
may ``include'' other header files, to incorporate additional
functionality.
\begin{verbatim}
/* my_code.h */
#ifndef MY_CODE
#define MY_CODE
#include <cmath>   // standard library math header
#include "epix.h"  // ePiX header
using ePiX::P;

namespace Mine {   // to avoid name conflicts
  // functions for special relativity
  double lorentz_norm(const P&);
  bool   spacelike(const P&);
} // end of namespace
#endif /* MY_CODE */
\end{verbatim}
This file exhibits two ``safety features''. The three \code{MY\_CODE}
lines prevent the file from being included multiple times. In a file
of this size, inclusion protection is overkill, but as your code base
grows and the number of header files increases, this protection is
essential. Second, the header introduces a ``Mine'' namespace. Inside
this namespace, two functions are declared as prototypes, giving the
function's return type, name, and argument type(s). A header file
should be commented fairly liberally, so that a year or two from now
you'll be able to decipher the file's contents.  For a longer file,
version and contact information, an overall comment describing the
file's features, and license information are appropriate.

Next, the corresponding source file; definitions are also placed into
the namespace, and must match their prototypes from the header file
exactly.
\begin{verbatim}
/* my_code.cc */
#include "my_code.h"
using namespace ePiX;

namespace Mine {
  double lorentz_norm(const P& arg)
  {
    double x(arg.x1()), y(arg.x2()), z(arg.x3()); // extract coords
    return (y-x)*(y+x) + z*z; // -x^2 + y^2 + z^2
  }
  bool spacelike(const P& arg)
  {
    return (lorentz_norm(arg) > 0); // true if inequality is
  }
} // end of namespace
\end{verbatim}
Copies of these files are included with the source code so you can
experiment with them. Next, the source file must be ``compiled'',
``archived'', and ``indexed''. In the commands below, the percent sign
is the prompt.
\begin{verbatim}
% g++ -c my_code.cc
% ar -ru libcustom.a my_code.o
% ranlib libcustom.a
\end{verbatim}
Please see your system documentation for details on command options
and what each step does. For linking (below), the name of the library
file must begin ``lib'' and have the extension~\filename{.a}. Once these
steps are successfully completed, put the library \filename{libcustom.a}
and header file \filename{my\_code.h} in your project directory. You're
ready to use the code in an \ePiX\ figure.


\subsection{Runtime Linking}

The script~\epix\ allows input files to be linked with external
libraries at run time, when the input file is compiled into a
temporary executable.

\epix\ recognizes command line options and passes them verbatim to the
compiler. The most commonly used options are those of the form
\begin{verbatim}
   -I<include>     -L<libdir>     -l<lib>
\end{verbatim}
For example, to link \filename{figure.xp} against
\filename{mylibs/libcustom.a}, run the command
\begin{verbatim}
   epix -Lmylibs -lcustom figure
\end{verbatim}
The options \code{-I. -L.} tell the compiler to look in the current
directory for header and library files.  Compiler options may appear
in any order, but must come before the name of the input file(s).

Compiler options may be placed in the configuration file
\filename{\$HOME/.epixrc}, with syntax as above. A line in the config
file that contains a pound sign~(\code{\#}) is a comment, no matter
where in the line the~\code{\#} appears. If any non-comment line
fails to start with a dash, the rest of the file is silently
discarded.  Command-line options are read before the config file.


\subsection{Using Multiple Versions}

The script \filename{epix}~links by default against the \code{C}~math
library \code{libm.a}~and the \ePiX\ library \code{libepix.a}.  The
command option \code{--no-defaults} clears the header and include
paths and removes \code{libepix.a} from the link list. The script may
therefore be used with multiple versions of \ePiX, a potentially
useful feature if you regularly need to compile old source files, or
simply prefer the syntax of an older version.

To install and use (say) Version~1.0.0, build the package according to
its \filename{INSTALL} instructions, but \emph{do not use the makefile
  to install}. Instead, manually install the header and library only,
using their version number:
\begin{verbatim}
# install -m 644 epix.h /usr/local/include/epix-1.0.h
# install -m 644 libepix.a /usr/local/lib/libepix-1.0.a
\end{verbatim}
A non-system directory may be used instead of \filename{/usr/local}.
To use the old version, a source file must \code{include} the
appropriate header file (which is identified by its version number).
To compile, issue a command such as
\begin{verbatim}
epix --no-defaults -I/usr/local -L/usr/local -lepix-1.0 file.xp
\end{verbatim}


\section{Programmer's Guide}
\label{section-programmers}

This section briefly surveys \ePiX's implementation, and is intended
for (potential) programmers. The source code is divided into modules
with small, well-defined responsibilities, but the user interface is
mostly compatible with the syntax of Version~1.0. These constraints
demanded a degree of implementation hiding. For example, the
user-visible classes defined in \filename{Color.h}, \filename{path.h},
and~\filename{screen.h} contain only a pointer to the implementation
class, and style data is hidden behind global commands.

Functionally, the code consists of the user interface; implementation
class\-es comprising drawing attributes, spatial objects, screens and
representations of their elements, and output; and miscellaneous
utility functions. The headers in each group, and their contents, are
described in Section~\ref{section-interface}.

The user interface headers are assembled into a single file,
\filename{epix.h}, and installed in \filename{/usr/local/include} by
default. In normal use, the shell scripts read only the user interface
header.  The individual headers, including the components of
\filename{epix.h}, are installed in
\filename{/usr/local/include/epix}. These are provided for authors of
external libraries, who may need access to implementation details.


\subsection{External Packages}

\ePiX\ harnesses the computational power of~\CXX\ to the typographical
capabilities of \LaTeX. Consequently, \ePiX\ should be viewed in part
as a framework for expressing numerical data visually.

In the course of your work, you may develop specialized code filling a
gap in \ePiX's functionality. If your code seems likely to be of
interest to other users, please consider bundling it as an external
package and notifying the \ePiX\ community so your work can be linked
from the project pages and distributed to interested users.

There are no formal requirements for external packages, but in the
interest of uniformity contributed code should follow the GNU Coding
Standards~\cite{GNUCoding}. At a minimum, an external package should
build with the standard \code{./configure; make; make install}
commands, and the \code{configure} script should accept an option
\code{--with-epix} for the user to specify a non-default \ePiX\
install directory.

If an external package builds a static library, it should provide a
single header containing all the package's entry points, and enclose
its interface in a namespace. Entry points should not collide with
\ePiX\ functions. Naming the package ``\filename{epix-<...>}'' is a
good idea, but not essential. For example, a package providing textual
nodes and diagram layout might be named \filename{epix-nodes}.  Input
files would use the package with the lines
\begin{verbatim}
  #include "epix-nodes.h"
  using namespace ePiX-nodes;
\end{verbatim}
and be compiled with
\begin{verbatim}
  epix -lepix-nodes <file>
\end{verbatim}

The user and internal interfaces of \ePiX-1.2 are not likely to
change. Still, it's prudent to rely only on the user interface in
contributed code whenever possible. Doing so also simplifies your work
as an author; your library can simply \code{include} the user header
file, and deal only with high-level objects and drawing attributes.


\subsection{User Interface}
\label{section-interface}

\index{Header file listing|(}
These files (in order) comprise the global header \filename{epix.h}.

\Header{enums.h} Marker, alignment, Riemann integral, and vector field
types.

\Header{length.h} Physical lengths, conceptually a number and
two-letter \LaTeX\ length unit.

\Header{interval.h} Closed, open, and half-open interval ranges for
data culling.

\Header{triples.h} The \code{P} class.

\Header{Complex.h} The \code{Complex} class.

\Header{functions.h} Angle-sensitive trig functions, miscellaneous
utility functions, the \code{Deriv}~and \code{Integral} classes.

\Header{pairs.h} Screen locations and displacements, with complex
arithmetic operations.

\Header{Color.h} The Color class interface, named
primaries and constructors.

\Header{state.h} Angle mode; clipping and cropping; label styles;
filling; arrow head style; dot and tick sizes, dash length; line
style; line and base pen attributes; color-setting commands.

\Header{frame.h} Orthonormal bases.

\Header{domain.h} Coordinate boxes for function plotting.

\Header{camera.h} The camera.

\Header{screen.h} The screen class.

\Header{picture.h} Dimension-setting, offset, layout, decoration,
verbatim text, and output format commands.

\Header{markers.h} Point markers, axis labels, and coordinate axes.

\Header{axis.h} Coordinate axes and labels in various styles.

\Header{legend.h} Plot legends.

\Header{path.h} The path class.

\Header{curves.h} Polygons, arrows, ellipses, arcs, splines,
coordinate grids, and recursive fractal curves.

\Header{circle.h} The \code{Circle} shape object class.

\Header{plane.h} The \code{Plane} shape object class.

\Header{segment.h} The \code{Segment} shape object class.

\Header{sphere.h} The \code{Sphere} shape object class.

\Header{intersections.h} Shape object intersection operators.

\Header{plots.h} Plotting commands.

\Header{surface.h} Shaded surface plots.

\Header{data\_mask.h} Helper class for data pruning.

\Header{data\_file.h} Class for storing and representing data.

\Header{data\_bins.h} Class for sorting and counting data.

\Header{geometry.h} Latitudes and longitudes; spherical plotting,
arcs, polygons, and polyhedra; hyperbolic arcs.

\Header{Sline.h} The \code{Sline} (spherical line) class.


\subsection{Implementation Classes}

\Header{Color\_Base.h} The Color implementation interface.

\Header{Color\_CMY.h} The CMY color model.

\Header{Color\_CMYK.h} The CMYK color model.

\Header{Color\_Gray.h} The Gray color model.

\Header{Color\_Neutral.h} Each color class has a ``Neutral'' member
that converts colors to that model by filtering. This file defines the
unique ``model-less'' Neutral color for which filtering performs no
action.

\Header{Color\_RGB.h} The RGB color model.

\Header{Color\_Sep.h} Classes for CMYK separation.

\Header{active\_screen.h} Simple manipulator for the active screen.

\Header{picture\_data.h} Picture implementation: two \code{screen}s
(representing the \code{canvas}~and the output page), true dimensions
and offsets, pointer to output format, list of colors, and lists of
verbatim text to write before and after printing the \code{picture}
environment in the output file. For simplicity, the \code{screen}s and
dimensions are public; encapsulation from the user results from
``hiding'' this header.


\subsubsection*{Style Attributes}

Declaration-style attributes are maintained with functions returning
static references: \code{the\_angle\_style()},
\code{the\_arrowhead\_style()}, \code{the\_label\_style()},
\code{the\_mark\_size()}, \code{the\_paint\_style()}, and
\code{the\_path\_style()}. Each function is declared in the
analogously-named header.

\Header{angle\_units.h} Angular modes: \code{radians},
\code{revolutions}, and \code{degrees}.

\Header{arrow\_style.h} Arrow head style data: width, ratio, and
inset.

\Header{label\_style.h} Text object style: Label and mask colors,
padding, border color and width, alignment, font size, font face, and
angle.

\Header{marker\_style.h} Dot and tick sizes.

\Header{paint\_style.h} Line and base pens, fill color.

\Header{path\_style.h} Solid, dashed, dotted lines.

\Header{pen\_data.h} The pen class.


\subsubsection*{Objects}

\Header{arrow\_data.h} Arrow representation.

\Header{facet.h} Shaded surface elements.

\Header{label\_data.h} Text (label and marker) objects.

\Header{legend\_item.h} Items for \code{legend}s.

\Header{path\_data.h} Path implementation.

\Header{spline.h} The natural spline class.

\Header{spline\_data.h} Templates for quadratic and cubic splines.


\subsubsection*{Screen Representation}

Elements in a \code{screen} are represented polymorphically as
``\code{tile}s'', of which there are six types: \code{glyph} (textual
elements), \code{pen\_arrow} (arrows), \code{pen\_fill} (filled
regions), \code{pen\_line} (path-like elements), \code{verbatim}
(raw text), and \code{legend\_tile} (legends).

Border and background shape are dictated by the \code{screen\_mask}
class. The cropping algorithm assumes the contour of a screen mask is
convex.

\Header{affine.h} Affine maps.

\Header{cropping.h} Screen mask cropping.

\Header{glyph.h} Markers and labels.

\Header{legend\_tile.h} Screen representation of a \code{legend}.

\Header{mask\_diamond.h} Screen diamond mask.

\Header{mask\_ellipse.h} Screen elliptical mask.

\Header{mask\_rectangle.h} Screen rectangular mask.

\Header{pen\_arrow.h} Arrows.

\Header{pen\_fill.h} Filled regions.

\Header{pen\_line.h} Paths.

\Header{screen\_data.h} Screen implementation class.

\Header{screen\_mask.h} Screen mask interface.

\Header{tile.h} Screen element interface.

\Header{verbatim.h} Text in output stream.


\subsubsection*{Output}

Output is divided into a couple of high-level operations and several
``atomic'' low-level operations.  To create a new output
format, one need only implement the \code{format} interface for the
desired file type. Paths and filled regions may be implemented however
the output type dictates.

\Header{eepic.h} eepic macros.

\Header{fmt\_template.h} ``skeleton'' header for new output formats.

\Header{format.h} The output interface.

\Header{pst.h} PSTricks macros.

\Header{tikz.h} tikz macros.


\subsubsection*{Utilities}

\Header{Color\_Utils.h} Functions for setting color channel densities.

\Header{clipping.h} The clip box.

\Header{constants.h} Global constants: Line widths; dot, tick, and
dash sizes; arrowhead parameters; miscellaneous numerical constants
and internal parameters.

\Header{crop\_algorithms.h} Path and loop clipping templates.

\Header{deriv.h} Finite difference template.

\Header{edge\_data.h} Path element representation template.

\Header{errors.h} Warning and error messages.

\Header{frac.h} Rational numbers and operators.

\Header{halfspace.h} Halfspace cutting.

\Header{hatching.h} Filling regions in eepic.

\Header{lens.h} Camera lenses.

\Header{map.h} Wrappers for templated plotting.

\Header{plot\_algorithms.h} Plotting templates.

\Header{screen\_crop.h} Crop paths and loops.

\Header{utils.h} Truncation, date and time, line breaking.
\index{Header file listing|)}


\appendix
\texinfoappendix{A}{Software Freedom}
\index{Free software|(} 

Academics in general, and mathematicians in particular, depend on free
exchange of information. We prove theorems or establish experimental
results, write up formal accounts, place preprints on public file
servers, and submit papers to peer-reviewed journals. If accepted, the
results---data, techniques, methods of reasoning, citations, and
conclusions---are published in print and become part of the public
record, governed by copyright law. Libraries purchase journal
subscriptions, but researchers and scholars may use ideas from the
literature merely by giving appropriate citations in their own
work. ``Theft'' arises from false claims of authorship.

Carried over to software, the academic process would guarantee rights
similar to those provided by the GNU General Public License (GPL):

\begin{itemize}

\item (GPL~0) To run a program for any purpose.

\item (GPL~1) To study how the program works, and adapt it to your
  needs.

\item (GPL~2) To redistribute copies of the program.

\item (GPL~3) To improve the program, and release improvements to the
  public, so that the whole community benefits.

\end{itemize}

In reality, attitudes toward software differ markedly. Most academics
work on a proprietary platform, use proprietary software for research
and teaching, and share information with colleagues and students in
proprietary, even obfuscated, data formats.  Contrary to the academic
ethic, proprietary software licenses restrict access to information:
preventing users from learning how a program works internally
(``reverse engineering''), limiting the number of users who may run a
piece of software, and forbidding users from running (or sometimes
even installing) a purchased copy on multiple machines.

Restrictions on use aside, if one cannot examine a program's source
code, one cannot fully trust the output, any more than one can trust
(for purposes of scientific publication) results of a commercial
testing lab.

\texicomment
\noindent$\phantom{.}$\hfil\vbox{\hrule width 4in}\hfil
\texicommentend

Consider a hypothetical future world in which scholarly results are
disseminated like software. Instead of subscriptions, journals 
sell licenses granting readership to a specified number of
individuals. Photocopying an article for a class or research seminar
constitutes ``piracy'', though if the institution has purchased a
sufficiently large site license the teacher or speaker may bring the
physical volume to class and project the pages onto a screen.

Of course, reading an article is scarcely enlightening. Mathematics
papers contain only the statements of theorems. Merely opening the
journal binds the reader to a lengthy legal agreement, stating that
theorems be used only for specific purposes and threatening serious
legal consequences for attempting to discover the author's proofs.

The actions of a single student, employee, or faculty member can
expose an institution to a costly ``journal audit'' from the
Mathematical Society of America (MSA), with the institution
responsible for legal costs if the audit reveals license violations
anywhere in the organization.

Mathematicians who long for the Old Days when papers contained
proofs and were shared freely are dismissed as idealistic cranks
or labeled anti-business communists. Common knowledge asserts the
obvious superiority of proprietary journals, and the necessity of
licenses for keeping mathematicians gainfully employed.

\texicomment
\noindent$\phantom{.}$\hfil\vbox{\hrule width 4in}\hfil
\texicommentend

Back in our world, some vendors have attempted to placate opponents of
closed source with ``shared source'' licenses, under which one may
sign a non-disclosure agreement and subsequently examine source
code. In the future world analogy, a shared source agreement would
allow journal licensees to sign an NDA, then see the proofs of
theorems. Readers could thereby correct errors in proofs (benefitting
the publisher by improving the reliability of the journal), but would
be legally forbidden from using the ideas elsewhere (denying benefit
to other mathematicians).

Legally and conceptually codifying software as a commodity ignores a
fundamental reality: Like an idea or recipe, software can be copied
without loss of the original. The perception of ``theft'' by copying
arises from an artificial belief that software has an owner who must
be monetarily compensated each time a person acquires a copy.  The
nature of software does not enforce the ``sale'' model in the way
services and physical commodities do. It therefore seems
philosophically inappropriate to treat software as a commodity, and
perilous to conform the legal system to the enforcement of such a
model.

At its best, software enhances our productivity and
creativity. Sharing software, like sharing ideas, benefits a larger
number of people without detriment to existing users.  I hope this
modest program is, in conjunction with the much larger efforts of
others (especially Donald Knuth, Richard Stallman, and the many people
who have contributed to the authorship of \LaTeX\ and its packages),
useful to you in your mathematical work.

Please visit the Free Software Foundation, at
\URL{www.fsf.org}, to learn more about free software and how
you can contribute to its development and adoption.
\index{Free software|)}

\texinfoappendix{B}{Acknowledgments}

\ePiX\ is built on the work of many people, many of whom I am
unaware. The following people have contributed, sometimes unknowingly
but always generously:

\paragraph{Infrastructure} Donald Knuth, Conrad Kwok, Leslie Lamport,
Tim Morgan, Piet van~Oostrum, Sunil Podar, Richard Stallman, Till
Tantau, Herbert Vo\ss, Timothy van~Zandt

\paragraph{Enhancements} Jay Belanger, Robin Blume-Kohout, Julian
Gilbey, Marcus Hanwell, Yvon Henel, Svend Daug{\aa}rd Pedersen, Andrew
Sterian

\paragraph{Porting and packaging} Youn\`es Hafri (Crux), Julian Gilbey
(Debian); Tsuguru Kato (FreeBSD); Markus Dittrich, Danny van Dyk,
Christian Faulhammer, Olivier Fisette, Chris Gianelloni, Michael
Hanselmann, Marcus Hanwell, David Holm, Peter Johanson, Patrick
Kursawe, Tobias Scherbaum, Markus Ullmann (Gentoo); Guido Gonzato
(RPM); Rene Rebe (T2)

\paragraph{Debugging, advice, and other assistance} Maik Beckmann, Jay
Belanger, Felipe Paulo Guazzi Bergo, Karl Berry, Robin Blume-Kohout,
Aran Clauson, Patrick Cousot, Stephen Gibson, Julian Gilbey, Dov
Grobgeld, Bob Grover, Jim Hefferon, Jacques L'helgoualc'h, Yvon Henel,
Hartmut Henkel, Herng-Jeng Jou, Walter Kehowski, Paul Kornman, Kevin
McCormick, Ross Moore, Mike Protts, Thorsten Riess, Jean-Michel
Sarlat, Alan Sill, Neel Smith, Michael Somos, Andrew Sterian, Ryszard
Tanas, Ben Tillman, Kai Trukenmueller, Torbjorn Vik, Wenguang Wang,
Gabe Weaver, Mariusz Wodzicki

\texicomment
\begin{thebibliography}{9}
%
\bibitem{Cline} Marshall Cline, \emph{\CXX\ FAQ Lite}, \\
  \URL{http://www.parashift.com/c++-faq-lite/}
%
\bibitem{xcolor} Uwe Kern, \emph{Extending \LaTeX's color facilities:
  the \code{xcolor} package}, white paper, Jan.~21, 2007
%
\bibitem{KnR} Brian Kernighan and Dennis Ritchie, \emph{The \code{C}
    Programming Language}, Second Ed., Prentice-Hall Software Series,
    1988
%
\bibitem{GNUC} Sandra Loosemore, Richard M. Stallman, et.\ al.,
    \emph{The GNU \code{C}~Library Reference Manual}, GNU Press,
    2004
%
\bibitem{Prata} Stephen Prata, \emph{\CXX\ Primer Plus}, Sams, 2002
%
\bibitem{R} Keith Reckdahl, \emph{Using Imported Graphics in
  \LaTeX2e}, Version~2.0, white paper, Dec.~15, 1997
%
\bibitem{GNUCoding} Richard M. Stallman, et.\ al., \emph{The GNU
  Coding Standards}, \\
  \URL{http://www.gnu.org/prep/standards/}
%
\bibitem{S} Bjarne Stroustrup, \emph{The \CXX\ Programming Language},
  Special Ed., Addison-Wesley, 1997
%
\bibitem{Z} Timothy van~Zandt, \emph{PSTricks: PostScript Macros for
  Generic \TeX}, Version~0.93a, white paper, Mar.~12, 1993
%
\end{thebibliography}

%\clearpage

%\printindex[xfcns]

\printindex

\end{document}

\texicommentend
